\ifx\danslelivre\undefined
\documentclass[10pt]{../fichiers-configuration-latex/smfart}
\usepackage{ifthen}
\newboolean{english}
\newboolean{francais}
\setboolean{english}{false}
\setboolean{francais}{true}
\input{../fichiers-configuration-latex/paquets}
\selectlanguage{francais}
\input{../fichiers-configuration-latex/lemme-proposition-theoreme}
\input{../fichiers-configuration-latex/numerotation}
\input{../fichiers-configuration-latex/macros}

\renewcommand\theequation{\thesubsection.\textbf{\alph{equation}}}

\externaldocument{../14-fonctions-dimension/fonctions-dimension}
\externaldocument{../21-finitude-non-abelienne/finitude-non-abelienne}


\begin{document}

\begin{center} XX. Rigidité\\ Yves Laszlo et Alban Moreau \end{center}

\date{exposés oraux des 22 juin 2007 et 11 juin 2008, notes du 29 août 2008}

\setcounter{tocdepth}{2}
\tableofcontents

\else
\setboolean{english}{false}
\setboolean{francais}{true}
\selectlanguage{francais}
\chapter{Rigidité}
\addtocontents{toc}{par Yves Laszlo et Alban Moreau}
\begin{center} Yves Laszlo et Alban Moreau
\end{center}
\fi


\section{Introduction} Le but de cet exposé est de démontrer les deux résultats techniques
\ref{20-prop-rig} (comparaison des torseurs sur l'ouvert complémentaire $\Spec{A}-V(I)$ défini par
un couple hensélien non nécessairement noethérien et l'ouvert correspondant du complété formel le
long de $V(I)$) et \ref{20rigfort} (rigidité de la ramification). Ils permettront dans l'exposé
suivant de montrer l'énoncé de finitude suivant (\refext{XXI}{21-local}):

\begin{theoreme0}[Gabber]   Soit $A$ un anneau strictement local de dimension $2$.  On suppose que
$A$ est normal, excellent, et on note $X'=\Spec(A)-\{\mathfrak{m}_A\}$  son spectre épointé.  Alors,
pour tout groupe fini $G$,   l'ensemble $H^1(X';G)$ est fini.\end{theoreme0}

Ce résultat est la clef pour démontrer le résultat de finitude général suivant
(\refext{XXI}{21-groupes}) :

\begin{theoreme0}[Gabber]    Soit $f:Y \rightarrow X$ un morphisme de type fini   entre schémas
quasi-excellents.  Soit $\mathbb L$ un ensemble de nombres premiers inversibles sur $X$.  Pour tout
faisceau constructible de groupes $F$ sur $Y_\et$    de $\mathbb L$-torsion,  le faisceau $\R^1
f_*(F)$ sur $X_\et$ est constructible. \end{theoreme0}

Par des techniques d'ultrafiltres, chères aux théoriciens des modèles, on est ramené à étudier des
revêtements étales de spectres épointés d'anneaux \textit{non noethériens}, ce qui explique qu'on
soit contraint de démontrer les énoncés techniques hors de tout cadre noethérien.

\begin{remarque0}\label{precau-champ} Soit $X$ un schéma. On considérera des champs en groupoïdes
$\C$ sur $X_\et$ (on dira simplement champs).  En général, la catégorie fibrée $\C$ n'est pas
scindée de sorte que si $x,y$ sont deux objets de $\C(S)$ où $S\ra X$ est étale, il faut quelques
précautions pour parler du faisceau  $\underline{\Hom}(x,y)$ sur $S_\et$. Précisément, suivant
\cite[I.2.6.3.1]{CNA@Giraud}, on considère l'équivalence de catégories fibrées $\C\ra{\lf}\C$ entre
$\C$ et la catégorie libre ${\lf}\C$ engendrée par $\C$, catégorie libre qui elle est scindée. On
définit alors $$\underline{\Hom}(x,y)(S')=\Hom_{{\lf}C(S')}({\lf}x',{\lf}y')$$ où
${\lf}x',{\lf}y'$ sont les images inverses par le morphisme étale $S'\ra S$ de $\lf x,\lf y$ dans
$\lf\C(S')$. Bien entendu (\cite[I.2.6.3.2 (1)]{CNA@Giraud}, $\lf$ induit une bijection
$$\Hom_{\C(S)}(x,y)\iso H^0(S,\underline{\Hom}(x,y)).$$ Ces remarques justifient qu'on puisse
si besoin supposer  sans dommage que les champs que l'on considérera sont scindés.  \end{remarque0}

\section{Lemme de rigidité} Soit $(A,I)$ un couple hensélien (\egalong{IV}{4}{18.5.5}) non
nécessairement noethérien, \textbf{avec $I$ de type fini}\footnote{Cette hypothèse sera utilisée
pour comparer les gradués $I$-adiques de $A$ et de son complété $\chap{A}$ (\ac{III}{2}{12}{})}.
Soit $U$ un ouvert de $X=\Spec(A)$ contenant $\Spec(A)-V(I)$. On note $\chap{A}$ le complété\footnote{On dira
simplement complété pour séparé complété.} $I$-adique de $A$ et $\chap{U}$ l'image inverse de $U$
par le morphisme de complétion $\pi:\chap{X}=\Spec(\chap{A})\ra X$.
On suppose pour simplifier $U$ quasi-compact (cf. \ref{20-U-non-qc}).

\subsection{Énoncés}\label{20-sec-enonc} Rappelons \sga{4}{IX}{1.5} qu'un faisceau en groupes $\F$
sur $X$  est ind-fini si pour tout ouvert étale $u:\ U\rightarrow X$ avec $U$ quasi-compact, le
groupe $\F(u)$ est limite inductives filtrante de ses sous-groupes d'indice fini. On dit alors qu'un
champ en groupoïdes $\C$ sur $X$ est ind-fini si pour tout si pour tout ouvert étale $u:\
U\rightarrow X$ avec $U$ quasi-compact et tout $x_u\in\C(u)$, le faisceau en groupes
$\pi_1(\C,x_u)=\Aut_\C(x_u)$ est ind-fini.

Le but de cette de cette section est de démontrer le théorème de rigidité suivant.




\begin{theoreme2}[Théorème de rigidité de Gabber] \label{20-prop-rig} Soit $\F$ un faisceau
d'ensembles sur $U_\et$. Alors on a \begin{itemize} \item[i)] la flèche naturelle $H^0(U,\F)\ra
H^0(\chap{U},\pi^*\F)$ est bijective ; \item[ii)] si $\F$ est de plus un faisceau en groupes
ind-fini, la flèche naturelle $H^1(U,\F)\ra H^1(\chap{U},\pi^*\F)$ est bijective.  \end{itemize}
\end{theoreme2}

Les deux énoncés du théorème précédent sont conséquence du théorème suivant, apparemment plus fort,
forme champêtre du théorème de rigidité\footnote{Les champs (ind-finis) en groupoïdes discrets
s'identifient aux faisceaux d'ensembles : on dira parfois un \textit{champ discret}.}.


\begin{theoreme2}[Théorème de rigidité de Gabber, forme champêtre] \label{20-prop-rig-champ} Soit
$\C$ un champ en groupoïdes ind-fini sur $U_\et$. Alors, la flèche naturelle
$\gamma(\C):\Gamma(U,\C)\ra \Gamma(\chap{U},\pi^*\C)$ est une équivalence.\end{theoreme2}



\begin{remarque2} En fait, le théorème~de rigidité \ref{20-prop-rig} est \textit{a priori}
équivalent à la version champêtre \ref{20-prop-rig-champ}. C'est ce qui ressort par exemple de
l'énoncé~\ref{20-prop-red-champ}. Mais, formellement, on n'a pas besoin de démontrer cela à ce
stade.\end{remarque2}

\begin{remarque2}
\label{20-U-non-qc}
 Les résultats précédents sont également valables lorsque $U$ n'est
pas nécessairement quasi-compact. Cela résulte du fait que
la catégorie des sections d'un champ sur $U$ est
équivalente à la $2$-limite projective des sections
sur les ouverts quasi-compacts de $U$ contenant $\Spec(A)-V(I)$.
L'hypothèse de quasi-compacité est utilisée dans un argument
d'éclatement ci-dessous (cf. \ref{20-fin-preuvebis}).
\end{remarque2}


\subsection{Réduction au cas constant}

Le résultat est le suivant

\begin{proposition2}\label{20-reduc-const}Supposons que pour tout $U$ comme plus haut,
\begin{itemize} \item[i)] pour tout ensemble fini $F$, la flèche $H^0(U,F)\ra H^0( \chap{U},F)$ est
bijective. Alors, \ref{20-prop-rig} i) est vrai, c'est-à-dire le théorème de rigidité
\ref{20-prop-rig-champ} est vrai pour les champs discrets.  \item[ii)]  pour tout groupe fini $G$,
la flèche $\Tors(U,G)\ra\Tors(\chap{U},G)$ est une équivalence et \ref{20-prop-rig} i) est vrai.
Alors le théorème de rigidité \ref{20-prop-rig-champ} est vrai.  \end{itemize} \end{proposition2}
\begin{proof} D'après \sga{4}{XII}{prop.6.5} (resp. \ref{20-prop-red-champ}), il suffit pour prouver
\ref{20-prop-rig} i) (resp. \ref{20-prop-rig-champ}) de prouver que pour tout $U'\ra U$ fini et tout
ensemble fini $F$ (resp. groupe fini $G$), la flèche \begin{equation}\label{20-lemme-sga}
H^0(U',F)\ra H^0(\chap{U'},F)\textup{ (resp. } \Tors(U',G)\ra\Tors(\chap{U'},G)) \end{equation}

 est bijective (resp. une équivalence) où $\chap{U'}=\chap{U}\times_UU'$.

\begin{lemme2} Il existe un schéma schéma affine $\Spec(B)$ et un diagramme cartésien
$$\xymatrix{U'\ar[r]\ar[d]\ar@{}[rd]|{\square}&\Spec(B)\ar[d]\\U\ar[r]&\Spec(A)}$$ où  $B$ est fini
sur $A$. Le morphisme $U'\ra \Spec(B)$ s'identifie à l'immersion ouverte $U_B\hookrightarrow
\Spec(B)$. De plus, $U_B$ contient $\Spec(B)-V(IB)$.  \end{lemme2}

\begin{proof} Comme $U'\ra U$ est fini, il est projectif (\ega{II}{6.1.11}). Comme $U$ est
quasi-compact, l'immersion ouverte $U\hookrightarrow X$ est quasi-affine (\ega{II}{5.1.1}), donc
quasi-projective de sorte que le composé $f:U'\rightarrow U\ra X$ est quasi-projectif
(\ega{II}{5.3.4}). Comme $X=\Spec(A)$ est affine, $\OO_X$ est certainement ample (cf. la définition
ou \ega{II}{5.1.2}). Les hypothèses du  théorème principal de Zariski (\egalong{IV}{3}{8.12.8}) sont
donc vérifiées. Il existe donc $X'\ra X$ fini de sorte que $f$ se factorise en $U'\hookrightarrow
X'\ra X$ où $U'\hookrightarrow X'$ immersion ouverte et $X'\ra X$ fini. L'adhérence schématique de
$U'$ dans $X'$ est fermée dans $X'$ : elle s'écrit donc $\Spec(B)$ où $B$ est fini sur $A$. On a
donc un diagramme commutatif $$\xymatrix{U'\ar[r]\ar[rd]&U_B\ar[d]\ar@{^(->}[r]&\Spec(B)\ar[d]\\
&U\ar@{^(->}[r]&\Spec(A)}$$ où les flèches non horizontales sont finies. La flèche $U'\ra U_B$ est
donc propre. Comme c'est aussi une immersion ouverte d'image dense, c'est un isomorphisme.  L'ouvert
$U$ contenant $\Spec(A)-V(I)$, on déduit que $U'=U_B$ contient
$\Spec(B)-V(IB)=(\Spec(A)-V(I))_B$.\end{proof}

D'après le lemme, la flèche (\ref{20-lemme-sga}) s'identifie à
\begin{equation}\label{20-lemme-sga-1} H^0(U_B,F)\ra H^0(\chap{U}_B,F)(\textup{ resp. }
\Tors(U_B,G)\ra\Tors(\chap{U}_B,G)) \end{equation}

 (où $?_B$ est l'extension des scalaires du $A$-schéma $?$ à $\Spec(B)$).  Il s'agit donc de montrer
que (\ref{20-lemme-sga-1}) est bijectif (resp. une équivalence).

Par définition, on a $$\chap{U}_B=\pi_C^{-1}(U)$$ où $\pi_C$ est la projection naturelle $$\pi_C:\
\Spec(C)\ra\Spec(A),\text{ avec } C=\chap{A}\otimes_AB.$$

Dans le cas noethérien, $C$ est le complété $IB$-adique $\chap{B}=C $ de $B$ ce qui prouve la
proposition dans ce cas -- appliquer l'hypothèse \ref{20-reduc-const} i) à $\F$ constant de valeur
$F$ sur $U_B$ --. Dans le cas général, la flèche $C\ra\chap{B}$ n'est pas en général un
isomorphisme.

\begin{lemme2}\label{20-lemme-hens}Avec les notations précédentes, on a \begin{itemize} \item [i)]
Soit $(A_n,I_n)$ un système projectif de couples henséliens. Le couple
$(A_\infty,I_\infty)=(\varprojlim{A_n},\varprojlim{I_n})$ est hensélien.  \item [ii)] Le complété
$I$-adique de $A$ est hensélien.  \item [iii)] Les couples $(B,IB)$ et $(C,IC)$ sont henséliens et
ont même complété $I$-adique.  \end{itemize} \end{lemme2} \begin{proof} Soit $P$ un polynôme
unitaire de $A_\infty[x]$ et $\bar a\in A_\infty/I_\infty$ une racine simple (c'est-à-dire  telle
que $P'(\bar a)$ inversible dans $A_\infty/I_\infty$). L'image $\bar a_n$ de $\bar a$ dans $A_n/I_n$
est une racine simple de $P$. Elle se relève donc de façon unique en une racine $a_n\in A_n$ de $P$
d'après le lemme de Hensel. Comme $I_{n+1}$ s'envoie dans $I_n$, par unicité des relèvements,
l'image de $a_{n+1}$ dans $A_n$ est égale à $a_n$ de sorte que la suite $a=(a_n)\in A_\infty$ est le
relèvement cherché de $\bar a$ ce qui prouve i).

Puisque $A\ra A/I^n$ est notoirement entier, les couples $(A/I^n, IA/I^n)$ sont henséliens de sorte
que ii) découle de i).


Par associativité du produit tensoriel, le morphisme naturel $B/I^nB\rightarrow C/I^nC$ s'identifie
au morphisme naturel $A/I^nA\rightarrow\chap{A}/I^n\chap{A}$. Comme ce dernier est un isomorphisme
(\ac{III}{2}{2}{, prop. 15 et cor. 2}), $B$ et $C$ ont même complété $I$-adique. iii) suit alors de
ii) car un couple fini sur un hensélien est hensélien.  \end{proof}

On a donc $\widehat{U_B}=\widehat{U_C}$. D'après le lemme précédent, sous les hypothèses de
\ref{20-reduc-const} i) (resp. ii)), la flèche naturelle $$H^0(U_B,F)\ra H^0(\widehat{U_B},F)=
H^0(\widehat{U_C},F)\leftarrow H^0({U_C},F)=H^0(\chap{U}_B,F)$$

(resp. $$\Tors(U_B,G)\ra \Tors(\widehat{U_B},G)= \Tors(\widehat{U_C},G)\leftarrow
\Tors({U_C},G)=\Tors(\chap{U}_B,G))$$ est alors une bijection (resp. équivalence), ce qu'on
voulait.\end{proof}



\subsection{Réduction au cas strictement hensélien} Résumons les notations dans le diagramme
cartésien suivant

$$\xymatrix{
\chap{U}\ar[r]^{\pi}\ar@{^(->}[d]_{{\chap{j}}}\ar@{}[rd]|\square\ar[r]^{\pi}&U\ar@{^(->}[d]^{j}\\
\chap{X}\ar[r]^{\pi}& X }$$ avec $U$ quasi-compact contenant $\Spec(A)-V(I)$.  Montrons le résultat
suivant.

\begin{proposition2}\label{20-reduchensel} Supposons que pour tout $U$ comme plus haut,
\begin{itemize} \item[i)] pour tout ensemble fini $F$, la flèche $H^0(U,F)\ra H^0(\chap{U} ,F)$ est
bijective si $A$ est de plus strictement local. Alors, \ref{20-prop-rig} i) est vrai (que $A$ soit
strictement local ou non).  \item[ii)]  pour tout groupe fini $G$, la flèche
$\Tors(U,G)\ra\Tors(\chap{U},G)$ est une équivalence si $A$ est de plus strictement hensélien et
\ref{20-prop-rig} i) est vrai. Alors, \ref{20-prop-rig-champ} est vrai (que $A$ soit strictement
local ou non).  \end{itemize} \end{proposition2}


\begin{proof} Commençons par un lemme.

\begin{lemme2}\label{20-cbase1} Supposons que pour tout $U$ comme plus haut, \begin{itemize}
\item[i)] pour tout ensemble fini $F$, la flèche $H^0(U,F)\ra H^0(\chap{U},F)$ est bijective si $A$
est de plus strictement local. Alors, la flèche de changement de base $$\gamma:\pi^*j_*F\ra
{\chap{j}}_*\pi^*F={\chap{j}}_*F$$ est un isomorphisme (que $A$ soit strictement local ou non).
\item[ii)]  pour tout groupe fini $G$, la flèche $\Tors(U,G)\ra\Tors(\chap{U},G)$ est une
équivalence si $A$ est de plus strictement hensélien et \ref{20-prop-rig} i) est vrai. Alors, la
flèche de changement de base $$\gamma:\pi^*j_*\underline{\Tors}(U,G)\ra
{\chap{j}}_*\pi^*\underline{\Tors}(U,G)
={\chap{j}}_*\underline{\Tors}(\chap{U},G),$$
où l'égalité résulte de \cite[III.2.1.5.7]{CNA@Giraud},
est une équivalence (que $A$ soit
strictement local ou non).
\end{itemize} \end{lemme2}
\begin{proof}


Les formules $j^*j_*=\Id\text{ et }\chap{j}^*\chap{j}_*=\Id$ assurent qu'on a
$$\chap{j}^*\pi^*j_*=\pi^*j^*j_*=\pi^*=\chap{j}^*\chap{j}_*\pi^*$$ de sorte que l'image inverse sur
$\chap{U}$ de la flèche de changement de base

\begin{equation}\label{20-CTCU} \pi^*j_*\C\ra {\chap{j}}_*\pi^*\C \end{equation} est une équivalence
pour tout champ en groupoïdes $\C$.


Soit $\hat x$ un point géométrique de $\chap{X}$ d'image le point géométrique $x=\pi\circ\hat x$ de
$X$ et montrons que la fibre de la flèche de changement de base  (\ref{20-CTCU}) en $\hat x$ est une
équivalence. D'après ce qui précède, on peut supposer $\hat x\not\in\chap{U}$. En particulier, $x\in
V(I)$.

Soit $A^{\hs}$ (resp. $X_{(x)}$)  l'hensélisé strict de $A$ (resp. $X$) en $x$ et $\chap{A}^{\hs}$
(resp. $\chap{X}_{(\hat x)}$) celui de $\chap{A}$ (resp. $\chap{X}$) en $\hat x$. On a un diagramme
commutatif où les flèches sont les flèches de fonctorialité, complétion ou stricte hensélisation

$$\xymatrix{\widehat{\chap{X}_{(\hat x)}}\ar[r]\ar[d]&\chap{X}_{(\hat
x)}\ar[r]\ar[d]&\chap{X}\ar[d]\\ {\chap{X}_{(\hat x)}}\ar[r]&X_{(x)}\ar[r]&X}$$

On note alors

$$\xymatrix{\widehat{\chap{U}_{(\hat x)}}\ar[r]\ar[d]&\chap{U}_{(\hat
x)}\ar[r]\ar[d]&\chap{U}\ar[d]\\ {\chap{U}_{(\hat x)}}\ar[r]&U_{(x)}\ar[r]&U}$$ l'image inverse du
diagramme par l'immersion ouverte $U\rightarrow X$. En particulier, $U_{(x)}$ (resp.
$\chap{U}_{(x)}$) désigne l'image inverse de l'hensélisé strict $X_{(x)}$ (resp. $\chap{X}_{(x)}$)
de $X$ (resp. $\chap{X}$) en $x$ (resp. $\hat x$) par $j$ (resp. ${\chap{j}}$). Comme $U$ est
quasi-compact, il en est de même des ouverts $U_{(x)},\chap{U}_{(\hat x)}$ de
$X_{(x)},\chap{X}_{(\hat x)}$ .

Les morphismes $j,{\chap{j}}$ étant cohérents, dans le cas i), la fibre $\gamma_{\hat x}$
s'identifie à la flèche naturelle $$H^0(U_{(x)},F)\ra H^0(\chap{U}_{(\hat x)},F)$$ tandis que dans
le cas ii) elle s'identifie à $$\Tors(U_{(x)},G)\ra \Tors(\chap{U}_{(\hat x)},G).$$

  On déduit que les flèches naturelles $$H^0(U_{(x)},F)\ra H^0(\widehat{ U_{(x)}},F)\text{ et }
H^0(\chap{U}_{(\hat x)},F)\ra H^0(\widehat{\chap{U}_{(\hat{x})}},F)$$ sont bijectives dans le cas i)
et que les flèches $$\Tors(U_{(x)},G)\ra \Tors(\widehat{U_{(x)}},G)\text{ et }\Tors(\chap{U}_{(\hat
x)},G) \ra \Tors(\widehat{\chap{U}_{(\hat{x})}},G)$$ sont des équivalences dans le cas ii). Il
suffit donc de voir que la flèche naturelle \begin{equation}\label{20iso-complet}
\widehat{\chap{U}_{(\hat{x})}} \ra\widehat{ U_{(x)}} \end{equation} est bijective, ou encore que
$$A^{\hs}\text{ et }\chap{A}^{\hs}\text{ ont même $I$-complété.}$$ Puisque l'anneau local
$\chap{A}^{\hs}$ est hensélien, il est a fortiori $I$-hensélien (exercice). Utilisant
(\ref{20-lemme-hens}), on constate que le $I$-complété $\widehat{\chap{A}^{\hs}}$ est hensélien.
Comme son corps résiduel est celui de $A^{\hs}$, il est strictement hensélien. La flèche $\chap{A}
\ra\widehat{A^{\hs}}$ induit donc une flèche $\chap{A}^{\hs}\ra\widehat{A^{\hs}}$ et donc, par
$I$-complétion, une flèche $$\widehat{\chap{A}^{\hs}}\ra\widehat{A^{\hs}}.\leqno (*)$$ Par ailleurs,
la flèche de complétion $A\ra\chap{A}$ induit par hensélisation stricte puis complétion une flèche
$$\widehat{A^{\hs}}\ra\widehat{\chap{A}^{\hs}}.  \leqno (**)$$ Les flèches (*) et (**) sont inverses
l'une de l'autre, d'où le lemme.\end{proof}




On a le diagramme commutatif à carré cartésien

$$\xymatrix{
\chap{U}\ar[r]^{\pi}\ar@{^(->}[d]_{{\chap{j}}}\ar@{}[rd]|\square\ar[r]^{\pi}&U\ar@{^(->}[d]^{j}\\
\chap{X}\ar[r]^{\pi}& X\\ X_I=\Spec(A/I)\ar@{^(->}[u]\ar@{^(->}[ru]& }$$

Comme on l'a observé, les paires $(A,I)$ et $(\chap{A},I\chap{A})$ sont henséliennes.  La flèche
$H^0(\chap{X},\C)\ra H^0(X_I,\C_{|X_I})$ est donc une équivalence pour tout champ ind-fini $\C$ sur
$X_\et$ d'après \cite[théorème 1']{Gab94}.


On déduit d'une part $$H^0(U,F)=H^0(X,j_*F)=H^0(X_I,(j_*F)_{|X_I})$$ et, d'autre part
$$H^0(\chap{U},F)=H^0(\chap{X},{\chap{j}}_*F)\stackrel{\ref{20-cbase1}}{=}H^0(\chap{X},\pi^*
j_*F)=H^0(X_I,(\pi^* j_*F)_{|X_I})$$ ce dernier n'étant autre que $H^0(X_I,(j_*F)_{|X_I})$ (bien
entendu l'isomorphisme induit $$H^0(U,F)\iso H^0(\chap{U},F)$$ est la restriction).

De même, on a
$$H^0(U,{\underline{\Tors}}(U,G))=H^0(X,j_*{\underline{\Tors}}(U,G))=H^0(X_I,j_*{\underline{\Tors}}(U,G)_{|X_I}))$$
et, d'autre part

\begin{eqnarray*} H^0(\chap{U},{\underline{\Tors}}(\chap{U},G))&=&
H^0(\chap{X},{\chap{j}}_*{\underline{\Tors}}(\chap{U},G)) \\ &\stackrel{\ref{20-cbase1}}{=}&
H^0(\chap{X},\pi^* j_*{\underline{\Tors}}(U,G))\\ &=&H^0(X_I,\pi^*
j_*{\underline{\Tors}}(U,G)_{|X_I})) \end{eqnarray*}

ce dernier n'étant autre que $H^0(X_I,j_*{\underline{\Tors}}(U,G)_{|X_I}))$ (bien entendu
l'isomorphisme induisant $$H^0(U,{\underline{\Tors}}(U,G))\iso
H^0(\chap{U},{\underline{\Tors}}(\chap{U},G)).$$ Reste à invoquer \ref{20-reduc-const}.\end{proof}

\subsection{Fin de la preuve de \ref{20-prop-rig-champ}} D'après~\ref{20-reduchensel}, pour prouver
\ref{20-prop-rig-champ}, il suffit de prouver l'énoncé suivant

\begin{proposition2}\label{20-fin-preuve}Supposons $A$ strictement hensélien (et
$I\subset\text{rad}(A)$) et soit $U$ comme plus haut.  \begin{itemize} \item[i)] pour tout ensemble
fini $F$, la flèche $H^0(U,F)\ra H^0(\chap{U},F)$ est bijective.  \item[ii)]  pour tout groupe fini
$G$, la flèche $\Tors(U,G)\ra\Tors(\chap{U},G)$ est une équivalence.\end{itemize} \end{proposition2}

La formule $\pi^*\underline{\Tors}(U,G)=\underline{\Tors}(\chap{U},G)$ (\cite{CNA@Giraud},
III.2.1.5.7) permet de réécrire \ref{20-fin-preuve} sous la forme suivante

\begin{proposition2}\label{20-fin-preuvebis}Supposons $A$ strictement hensélien (et
$I\subset\text{rad}(A)$) et soit $U$ comme plus haut. Désignons par $\C$ le champ discret $F_U$ ou
bien $\underline{\Tors}(U,G)$. Alors, la flèche $H^0(U,\C)\ra H^0(\chap{U},\pi^*\C)$ est une
équivalence.  \end{proposition2}

\begin{proof} On va se ramener par éclatement au cas ou l'idéal $J$ définissant le complémentaire de
$U$ est principal .

Pour tout idéal ${\tilde I}$ d'un anneau ${\tilde A}$, on note $$\Ecl_{{\tilde I}}({\tilde
A})=\Proj(\oplus_{n\geq 0}\tilde I^n)$$ l'éclatement de ${\tilde I}$ dans $\Spec(\tilde A)$.
Si ${\tilde I}$ est de type fini, le morphisme structural $e:\Ecl_{{\tilde I}}({\tilde
A})\ra\Spec(\tilde A)$ est projectif, en particulier propre. Si $S$ est un $A$-schéma, on note
$\chap{S}^J$ le complété de $S$ le long de $J$.


On suppose donc $A$ strictement hensélien de corps résiduel $k$ et $\F=F_U$ comme plus haut. On a
déjà observé que $\chap{A}$ était aussi strictement hensélien. Il suit en particulier que l'ensemble
des sections globales de tout faisceau étale sur $X$ ou $\chap{X}$ s'identifie à sa fibre spéciale,
ce qu'on utilisera sans plus de précaution.

Comme $U$ est quasi-compact, il existe un idéal $J$ de type fini tel que $U=\Spec(A)-V(J)$. Comme
$U$ contient $\Spec{A}-V(I)$ et que $I$ est de type fini, on peut supposer $I\subset J$.  Soit
$$Y=\Ecl_J(A)\text{ et }Y'=\Ecl_J(\chap{A}).$$ (On aurait dû écrire $\Ecl_{J\chap{A}}(\chap{A})$ pour
$\Ecl_J(\chap{A})$).  Pour des raisons de cohérences, on notera simplement $X'$ le complété
$\chap{X}=\Spec(\chap{A})$ (resp. $U'$ sa restriction $\chap{U}=\pi^{-1}(U)$ à $U$).

\begin{sous-lemme2}\label{20-compa-comple}Soient $n,m$ des entiers $\geq 0$. Le morphisme de
complétion définit des isomorphismes $$A/I^mJ^n\simeq \chap{A}/I^mJ^n\chap{A}\text{ et }A/J^n\simeq
\chap{A}/J^n\chap{A}$$ induisant un isomorphisme \begin{equation*} J^n/I^mJ^n\simeq
J^n\chap{A}/I^mJ^n\chap{A}.  \end{equation*} \end{sous-lemme2}

\begin{proof}Comme $I$ est de type fini, le morphisme de complétion induit des isomorphismes
$$A/I^{m+n}\simeq \chap{A}/I^{m+n}\chap{A}\text{ et } A/I^{n}\simeq \chap{A}/I^{n}\chap{A}$$ d'après
\ac{III}{2}{12}{, cor. 2 de la prop. 16}. Mais comme $J$ contient $I$, on a $$I^{m+n}\subset
I^mJ^n\text{ et }I^n\subset J^n,$$ de sorte que les changements de base $$A/I^{m+n}\rightarrow
A/I^mJ^n\text{ et }A/I^{n}\rightarrow A/J^n$$ donnent alors des isomorphismes  $$A/I^mJ^n\simeq
\chap{A}/I^mJ^n\chap{A}\text{ et }A/J^n\simeq \chap{A}/J^n\chap{A}$$ qui donnent
\ref{20-compa-comple}.  \end{proof} La flèche naturelle $Y'\ra Y$ est donc un isomorphisme au dessus
de $\Spec(A/I)\subset X$ car elle est induite par le morphisme gradué $$\oplus J^nA/IJ^{n}A\ra
J^n\chap{A}/IJ^{n}\chap{A}$$ qui est un isomorphisme. On identifiera ces restrictions par la suite.
En particulier, le morphisme $p_s:Y'_s\ra Y_s$ entre fibres spéciales (c'est-à-dire  au dessus du
point fermé de $s\in \Spec(A/I)\subset X$)  est un isomorphisme grâce auquel nous les identifierons.
Regardons le solide commutatif

$$\xymatrix{ \chap{U}=U'\ar@{^(->}[r]^{j'}\ar[dd]^{p}&Y'\ar[dd]^p\ar[rrd]^{e}&Y_s\ar@{=}[dd]
\ar[rrd]\ar@{_(->}[l]_{i'}\ar@{}[ldd]| \square\\
&&&\chap{X}=X'\ar[dd]^\pi&\Spec(k)\ar@{_(->}[l]\ar@{=}[dd]\\
U\ar@{^(->}[r]^{j}&Y\ar[rrd]^-e&Y_s\ar[rrd]\ar@{_(->}[l]_i\\ &&&X&\Spec(k)\ar@{_(->}[l]}$$

Admettons pour un temps le résultat suivant.

\begin{lemme2}\label{20-commut-princ} Soit $\C=F_U$
(resp. $\C=\underline{\Tors}(U,G)$).  Alors, la flèche de changement de base $$\gamma :p^*j_{*}\C\ra
j'_{*}p^*\C$$ est bijective (resp. une équivalence).\end{lemme2}

Déduisons alors l'isomorphisme cherché $$H^0(U,{\C})\iso H^0(U',{\C})=H^0(U',p^*{\C})$$ grâce au
théorème de changement de base propre d'Artin-Grothendieck (\cite{CNA@Giraud} dans le cas noethérien
et théorème \ref{20-th-cgt-base-propre} dans le cas général) appliqué aux faces inférieure et
supérieure du diagramme précédent.  On a en effet un diagramme commutatif où toutes les flèches sont
les flèches naturelles (obtenues par adjonction)

$$\xymatrix{H^0(U,{\C})\ar@{=}[r]\ar[d]^\alpha&H^0(Y,j_{*}{\C})\ar[r]^-{b}\ar[d]^a
&H^0(Y_s,i^*j_{*}{\C})\ar[d]\ar@{=}[r]&H^0(Y_s,i'^*p^*j_*{\C})\ar[ld]^c\\
H^0(U',p^*{\C})\ar@{=}[r]&H^0(Y',j'_{*}p^*{\C})\ar[r]^-{d} &H^0(Y_s,i'^{*}j'_{*}p^*{\C})}$$

Les flèches $b,d$ dont bijectives grâce au théorème de changement de base propre
(\ref{20-th-cgt-base-propre}) tandis que $c$ une bijection grâce à (\ref{20-commut-princ}). Il suit
que $a$ et $\alpha$ sont des bijections.

\begin{proof}\textit{Preuve du lemme \ref{20-commut-princ}:} Soit $x'$ un point géométrique de $Y'$
d'image $x$ dans $Y$. On peut supposer $x'\in V(J)$. Soit $B$ l'hensélisé (strict) de $Y$ en $x$ et
$B'$ celui de $Y'$ en $x$. On doit étudier la flèche $$H^0(\Spec(B)-V(JB),{\C})\ra
H^0(\Spec(B')-V(JB'),{\C})\leqno (*).$$ Observons que par définition de l'éclatement, $JB$ (resp.
$JB'$) est un idéal principal engendré par un élément non diviseur de zéro et non inversible $t\in
B$ (resp. $t'\in B'$) (équation locale du diviseur exceptionnel). Par ailleurs, les couples $(B,JB)$
et $(B',JB')$ sont henséliens car $B,B'$ sont locaux henséliens (exercice). Les isomorphismes
$$J^nA/J^{n+m}A\iso J^n\chap{A}/J^{n+m}\chap{A},n,m\geq 0$$ assurent  que $B$ et $B'$ ont même
complété $J$-adique $\chap{B}=\chap{B}^J$.

On utilise alors les généralisations des résultats d'Elkik \cite{Elk73} -- et donc de
Ferrand-Raynaud pour le $\pi_0$ -- au cas principal non noethérien de \cite{ART@Gabber-Ramero}.
Précisément, le théorème 5.4.37 appliqué au $B[t^{-1}]$-groupoïde discret
$F_B=\Spec(B[t^{-1}])\times F$ assure qu'on a
$$H^0(\Spec(B[t^{-1}]),F)=\pi_0(F_B)=\pi_0(F_{\chap{B}})=H^0(\Spec(\chap{B}[t^{-1}]),F)$$ et de même
en remplaçant $B,t$ par $B',t'$. Comme $B$ et $B'$ ont même complété $J$-adique, on a donc
$$H^0(\Spec(B[t^{-1}]),F)=H^0(\Spec(B'[t'^{-1}]),F),$$ ce qu'on voulait. Dans le cas
$\C=\underline{\Tors}(U,G)$, on déduit du cas discret que (*) est pleinement fidèle. Soit alors
$\chap{P}$ un revêtement galoisien de groupe $G$ sur $\chap{U}=\Spec(\chap{B})-V(J\chap{B})$.
D'après le théorème 5.4.53 de  \cite{ART@Gabber-Ramero}, il provient d'un (unique) revêtement $P$ de
$U$. La pleine fidélité de (*) assure que le groupe d'automorphismes de $P$ est $G$. Dire que $P$
est galoisien de groupe $G$, c'est dire que la flèche canonique $$\phi: P\times G\ra P\times_UP$$
est un isomorphisme. On peut voir cette flèche comme un morphisme de revêtements étales de $U$.
Après image inverse sur $\chap{U}$, elle s'identifie à la flèche analogue $$\chap{P}\times G\ra
\chap{P}\times_{\chap{U}}\chap{P}$$ qui est un isomorphisme (de revêtements étales de $\chap{P}$
donc de revêtements étales de $\chap{U}$) par hypothèse. La pleine fidélité de (*) assure que $\phi$
est un isomorphisme de sorte que $P$ est bien galoisien de groupe $G$. On a donc obtenu que le
foncteur naturel entre les catégories de $G$-revêtements galoisiens sur $U$ et $\chap{U}$ sont
équivalentes. Il en est donc de même pour le foncteur les catégories de $G$-revêtements galoisiens
sur $U'$ et $\chap{U'}$. On conclut en se souvenant de l'égalité $\chap{U}=\chap{U'}$. \end{proof}

\end{proof}


\begin{remarques2}\label{20CTC} Le théorème~\ref{20-prop-rig-champ} entraîne immédiatement que la
flèche de changement de base $$\pi^*j_*\C\ra\chap{j}_*\pi^*\C$$ est une équivalence. En effet, on
l'a déjà vu sur $\chap{U}$ (\ref{20-CTCU}). Si $\hat{x}\not\in\chap{U}$, on a déjà observé dans la
preuve de \ref{20-cbase1} que $U_{(\bar x)}$ et $\chap{U}_{(\overline{\hat x})}$ avaient même
complété $I$-adique de sorte que deux applications de \ref{20-prop-rig-champ} assurent que la fibre
de $$\pi^*j_*\C\ra\chap{j}_*\pi^*\C$$ en $\overline{\hat x}$ est une
équivalence.\end{remarques2}



\section{Rigidité de la ramification}

\subsection{La condition $c_2$}

Rappelons (\egalong{IV}{4}{18.6.7}) que l'hensélisé $A^\h$ d'un anneau semi-local $A$ est le produit
des hensélisés des localisés de $A$ en ses idéaux maximaux. Pour tout anneau noethérien, on note
$A^\nu$ son normalisé, à savoir la clôture intégrale de $A$ dans l'anneau total $K(A)$ des fractions
de $A_\red$. Puisque $K(A)$ est le produit des $K(A/\p)$ où $\p$ décrit les points maximaux de
$\Spec(A)$, le normalisé de $A$ est le produit des normalisés des $A/\p$. Si $A^\nu$ n'est en
général pas noethérien (\cite[exemple 5 de l'appendice]{Nag62}), il est en revanche quasi-fini sur
$A$ (\cite[V.33.10]{Nag62}). En particulier, si $A$ est local, $A^\nu$ est semi-local. Si $A$ est
local noethérien, $A^\nu$ est semi-local de sorte que son hensélisé est bien défini. On a alors
(comparer avec \cite[43.20 et exercice 43.21]{Nag62})

\begin{lemme}\label{20hnu=nuh}Soit $A$ un anneau local noethérien.  \begin{itemize} \item La flèche
canonique $A^\h\ra (A^\nu)^\h$ induit un isomorphisme $(A^\h)^\nu\iso (A^\nu)^\h$.  \item Cette
bijection induit une bijection canonique $\p\mapsto\p^*$ entre les points maximaux $\p$ de
$\Spec(A^h)$ et les points fermés $\p^*$ de $\Spec(A^\nu)$  de telle sorte que les anneaux intègres
$(A^h/\p)^\nu$ et $(A^\nu_{\p^*})^h$  sont (canoniquement) isomorphes.  \end{itemize} \end{lemme}



\begin{proof}D'après \egalong{IV}{4}{18.6.8}, le morphisme canonique $A^\nu\otimes_AA^h\ra
(A^\nu)^\h$ est un isomorphisme. Le morphisme canonique  $A\ra A^\h$ étant ind-étale, il est normal.
D'après \egalong{IV}{2}{6.14.4}, le morphisme canonique $A^\h\ra A^\nu\otimes_AA^h$ identifie
$A^\nu\otimes_AA^h$ à la fermeture intégrale de $A^\h$ dans $A^\h\otimes_AK(A)$. Si maintenant,
$A\ra B$ est étale, la fibre au point maximal $\p\in\Spec(A)$ s'identifie à $\Spec(K(B))$. En
passant à la limite, on déduit l'égalité $A^\h\otimes_AK(A)=K(A^\h)$ de sorte que
$A^\nu\otimes_AA^h$ s'identifie à la fermeture intégrale de $A^\h$ dans $A^\h\otimes_AK(A)=K(A^\h)$
et donc $(A^\h)^\nu\iso A^\nu\otimes_AA^h$. La composition $$(A^\h)^\nu\iso A^\nu\otimes_AA^h\iso
(A^\nu)^\h$$ est l'isomorphisme annoncé. Pour le second point, on observe d'une part que le spectre
du normalisé de $A^\h$ est la somme disjointe des normalisés de ses composantes irréductibles
\begin{equation}\label{20dec-norm} \Spec((A^\h)^\nu)=\coprod_{\p\textup{\ point
maximal}}\Spec((A^\h/\p)^\nu), \end{equation} chaque fermé $\Spec((A^\h/\p)^\nu)$ étant intègre
(puisque local et normal) de sorte que \ref{20dec-norm} est une la décomposition en composantes
irréductibles de $\Spec((A^\h)^\nu)$.  D'autre part, par définition de l'hensélisé d'un anneau
semi-local, on a \begin{equation}\label{20dec-hens} \Spec((A^\nu)^\h)=\coprod_{\p^*\textup{\ point
fermé}}\Spec((A^\nu_{\p^*})^\h).  \end{equation} Or, $(A^\nu_{\p^*})^\h$ est local et normal (comme
$A^\nu_{\p^*}$), donc intègre, prouvant que \ref{20dec-hens} est la décomposition en composantes
irréductibles de $\Spec((A^\nu)^\h)$. Le lemme suit.  \end{proof}

\begin{proposition}\label{20propc2} Soit $Z$ un sous-schéma fermé d'un schéma noethérien $X$.  Les
conditions suivantes sont équivalentes : \begin{itemize} \item [(i)] Soit $p:X^\nu\ra X$ le
morphisme de normalisation. Alors, $p^{-1}(Z)$ est de codimension $\geq 2$ dans $X^\nu$.  \item
[(ii)] Pour tout $z\in Z$, toutes les composantes irréductibles de $\Spec(\OO_{X,z}^\h)$ sont de
dimension $\geq 2$.  \item [(iibis)] Pour tout $z\in Z$, toutes les composantes irréductibles de
$\Spec(\OO_{X,z}^\hs)$ sont de dimension $\geq 2$.  \item [(iii)] Pour tout $z\in Z$, toutes les
composantes irréductibles de $\Spec(\widehat{\OO_{X,z}})$ sont de dimension $\geq 2$.

\end{itemize} \end{proposition}

\begin{proof}Notons $A=\OO_{X,z}$ pour $z\in Z$. Notons d'abord que le morphisme $A^\h\ra A^\hs$ est
injectif, entier et fidèlement plat. Ceci prouve que le morphisme $h:\Spec(A^\hs)\ra\Spec(A)$
vérifie $\dim(\overline{h(x)})=\dim(\overline{x})$ et induit une surjection au niveau des points
maximaux, ce qui prouve l'équivalence de (\textit{ii}) et (\textit{iibis}).

Un anneau intègre et son normalisé ainsi qu'un anneau local et son hensélisé, ont même dimension.
Conservant les notations de \ref{20hnu=nuh}, on a donc $$\dim A^h/\p=\dim A^\nu_{\p^*}.$$  Or, dire
$\codim p^{-1}(Z)\geq 2$, c'est dire $\dim A^\nu_{\p^*}\geq 2$ lorsque $\p^*$ décrit les points
fermés de  $$\Spec(A^\nu)=p^{-1}(\Spec(\OO_{X,z}))$$ lorsque $z$ décrit $Z$. Ceci revient donc à dire
que toutes les composantes irréductibles $\Spec(A^h/\p)$ de $\Spec(A^h)$ sont de dimension $\geq 2$
prouvant l'équivalence de (\textit{i}) et (\textit{ii}).

Pour montrer l'équivalence de (\textit{i}) et (\textit{iii}), on peut supposer que $X=\Spec(A)$ est
local hensélien et que $Z$ est réduit à son point fermé.

Prouvons d'abord que (\textit{iii}) implique (\textit{ii}). Soit $Y$ une composante irréductible de
$X$. Le morphisme de complétion $c:\chap{X}\ra X$ étant  fidèlement plat, $\chap{Y}=c^{-1}(Y)$ est
une réunion de composantes irréductibles de $\chap{X}$ de sorte qu'on a $\dim(\chap{Y})\geq 2$.
Comme $Y$ est local noethérien, on a $\dim(Y)=\dim(\chap{Y})\geq 2$.

Prouvons la réciproque. Quitte à se restreindre à une composante irréductible (réduite), on peut
supposer $X$ intègre de dimension $\geq 2$. Soit $\hat x$ (resp. $x$) le point fermé de $\chap{X}$
(resp. $X$) (ce n'est pas une composante irréductible de $\chap{X}$ qui est de dimension $\geq 2$).
Si une des composantes de $\chap{X}$ était de dimension $\leq 1$, elle serait de dimension $1$ (car
$\{\hat x\}$ n'est pas une composante) et donc son point générique serait un point isolé de
$\chap{X}-\{\hat x\}$ de sorte que $\chap{X}-\{\hat x\}$ serait disconnexe (étant de dimension $\geq
2$). Or, d'après \cite[corollaire 4.4]{Fer-Ray70}, la flèche $$\pi_0(\chap{X}-\{\hat
x\})=\pi_0(c^{-1}(X-\{x\}))\ra\pi_0(X-\{x\})$$ est bijective. Or, comme $X$ est intègre de dimension
$\geq 2$, l'ouvert  $X-\{x\}$ est intègre donc connexe. \end{proof}

\begin{definition}\label{20defc2} Avec les notations de \ref{20propc2}, si $Z$ vérifie les
conditions équivalentes de  \ref{20propc2}, on dit que $Z$ est $c_2$ dans $X$.  \end{definition}

\begin{remarque}\label{20remc2} Si $X$ est intègre et excellent, $Z$ est $c_2$ si et seulement si
$X-Z$ contient tous les points de codimension $\leq 1$. En effet, comme le morphisme de
normalisation est fini et $X, X^\nu$  caténaires, on a $\dim\OO_{X^\nu,z^\nu}=\dim\OO_{X,p(z^\nu)}$
pour tout $z^\nu\in p^{-1}(Z)$.\end{remarque}

\begin{proposition}\label{20platc2} Soit $f:X'\ra X$ un morphisme plat de schémas noethériens et $Z$
un fermé de $X$. Alors, si $Z$ est $c_2$ dans $X$, son image inverse $Z'=f^{-1}(Z)$ est $c_2$ dans
$X'$. En particulier, la condition $c_2$ est invariante par localisation Zariski ou
étale.\end{proposition}

\begin{proof}Soit $z'\in Z'$ d'image $z=f(z')\in Z$. On suppose donc (\ref{20propc2}) que toutes les
composantes de $A=\widehat{\OO_{X,z}}$ sont de dimension $\geq 2$ et on veut prouver que toutes les
composantes de $B=\widehat{\OO_{X',z'}}$ sont de dimension $\geq 2$. On peut donc supposer que $f$
est morphisme local de schémas noethériens, locaux et complets. Comme $f$ est plat, toute composante
de $X'$  domine une composante $X_0$ de $X$ et est une composante de $f^{-1}(X_0)$. On peut donc
supposer $X$ intègre de dimension $>1$, de point fermé $z$. D'après \sga{2}{VIII}{2.3}, le
$A$-module $\OO(X-z)$ est de type fini. Comme $B$ est plat sur $A$, on déduit que
$B\otimes_A\OO(X-z)=\OO(X'-f^{-1}(z)$ est de type fini sur $B$. Comme $B$ est noethérien, le sous
$B$-module $\OO(X'-z')$ de $\OO(X'-f^{-1}(z))$ est de type fini. Mais si une des composantes $X'_0$ de
$X'$ était de dimension $1$, le complémentaire $X'_0-z'$ serait réduit au point générique $\eta$ de
$X'_0$ qui serait isolé dans $X'-z'$ isolé de sorte que $\OO(X'_0-z')$ serait un sous $B$-module de
$\OO(X'-z')$, donc de type fini ($B$ est noethérien). A fortiori, $\OO(X'_0-z')$ serait de type fini
comme $\OO(X'_0)$-module, ce qui contredit \sga{2}{VIII}{2.3} puisque $X'_0$ est de dimension
$1$.\end{proof}

\section{Théorème de rigidité de la ramification I : forme faible } Nous allons commencer par
démontrer une variante du changement de base lisse qui est cruciale dans la preuve du théorème de
rigidité~\ref{20rigfaible}.  \subsection{Variante du théorème de changement de base lisse}


Soit $G$ un groupe fini. On va démontrer une variante du théorème de changement de base lisse
\sga{4}{XVI}{1.2} pour les faisceaux de $G$-torseurs sans hypothèse sur le cardinal de $G$, mais en
se restreignant au cas d'immersions ouvertes.

\begin{theoreme2}[Gabber]\label{20-cbl}Considérons un diagramme cartésien
$$\xymatrix{U'\ar@{^(->}[r]^{j'}\ar[d]\ar@{}[rd]|\square&X'\ar[d]^p\\ U\ar@{^(->}[r]^{j}&X}$$
Supposons $X$ excellent normal, $p:X'\ra X$ lisse et $j:U\ra X$  immersion ouverte telle que $U$
contient tous les points de codimension $\leq 1$. Alors, le morphisme de changement de base
$\Phi:p^*j_*\tors(U,G)\ra j'_*\tors(U',G)$ est une équivalence.  \end{theoreme2}

\begin{proof} D'après le théorème de changement de base lisse  pour les faisceaux d'ensembles
\sga{4}{XVI}{1.2}, $\Phi$ est pleinement fidèle. Il suffit de prouver l'essentielle surjectivité.
Soit $x'$ un point géométrique de $X'$ d'image $x=p(x')$. Passant aux fibres, on est ramené à
prouver que la flèche d'image inverse des torseurs $H^1(U_{(x)},G)\ra H^1(U'_{(x')},G)$ est
bijective, avec de plus $x'$ \emph{fermé} dans sa fibre \sga{4}{VIII}{3.13 b)}. La stricte
hensélisation préserve la normalité et la codimension (platitude).  Les propriétés de permanence des
anneaux excellents (cf. \refext{I}{1-henselisation-anneaux-excellents})
assurent donc qu'on peut supposer $X=\Spec(A), X'=\Spec(A')$
avec $A=\OO_{X,x}^\hs, A'=\OO_{X',x'}^\hs$ strictement locaux, normaux et excellents. Comme $p$ est
lisse, le choix de coordonnées locales $t_1,\cdots,t_n$ de $X'$ en $x'$ définit un $A$-isomorphisme
$A\{t_1,\cdots,t_n\}\iso A'$ où comme d'habitude $A\{t_1,\cdots,t_n\}$ désigne l'hensélisé strict de
$A[t_1,\cdots,t_n]$ à l'origine. Une récurrence évidente permet de supposer $n=1$. On s'est ramené à
la situation $$\xymatrix{{U'}\ar@{}[rd]|\square\ar[d]^p\ar@{^(->}[r]&X'\ar[d]^p\\
U\ar@{^(->}[r]^-{j}\ar@/^1pc/[u]_\sigma&X\ar@/_1pc/[u]_\sigma}$$ avec $A$ strictement local, normal
et excellent et $\sigma$ la section de $p$ définie par l'immersion fermée d'équation $t=0$. Comme
$X,X'$ sont locaux et normaux, ils sont intègres. Les ouverts non vides de $X,X'$ sont donc intègres
et donc connexes. Le composé $$\pi_1(U)\xrightarrow{\sigma_*}\pi_1({U'})\xrightarrow{p_*}\pi_1(U)$$
étant l'identité, il suffit de prouver que $\sigma_*$ est surjectif. Soit alors ${V'}$ un revêtement
étale connexe de ${U'}$. On doit prouver que sa restriction $V\ra U$ au fermé
$U\stackrel{\sigma}{\hookrightarrow} U'$ d'équation $t=0$ est connexe.

Comme ${X'}$ est excellent, la clôture intégrale ${Y'}$ de ${X'}$ dans ${V'}$  est finie sur ${X'}$,
normale et intègre (comme $X'$). Comme $X'$ est hensélien, il en de même de $Y'$ qui est donc une
union disjointe de ses composants locaux. Comme $Y'$ est intègre, $Y'$ est local.  Soit $D\subset
{Y'}$ le diviseur de Cartier d'équation $t=0$ :  $D$ est connexe, puisque fermé dans un schéma
local.


On a donc un diagramme commutatif à carrés cartésiens et où les flèches verticales sont finies (et
dominantes).

$$\xymatrix{V'\ar@{^(->}[r]\ar[d]&Y'\ar[d]&D\ar@{_(->}[l]\ar[d]&V\ar@{_(->}[l]\ar[d]\\
U'\ar@{^(->}[r]&X'&X\ar@{_(->}[l]_-\sigma&U\ar@{_(->}[l]}$$

Soit $x'$ un point de $D-V$, d'image $x$ dans $X-U$. Comme $D\ra X$ est fini, on a
$\dim\overline{\{x'\}}=\dim\overline{\{x\}}$ et $\dim(D)=\dim(X)$. Comme $X$ et $X'$ sont caténaires
(ils sont même excellents), on en déduit l'égalité $\dim\OO_{D,x'}=\dim\OO_{X,x}$ ce qui assure que
l'ouvert $V$ dans $D$ contient tous les points de codimension $1$ dans $D$ (de même que  le
complémentaire de $U$ dans $X$ contient tous les points de codimension $1$ dans $X$). D'après le
lemme \refext{XXI}{21-ref-deg-connexe} appliqué au diviseur de Cartier connexe du schéma normal,
excellent  ${Y'}$, le schéma $V$ est connexe.\end{proof}

\subsection{Énoncé et réductions}

\begin{theoreme2}[Rigidité de la ramification] \label{20rigfaible} Soient $X,X'$ des schémas
noethériens, $Z\subset X$ un sous-schéma fermé, $U\stackrel{{\tilde j}}{\hookrightarrow}X$ l'ouvert
complémentaire et $X'\xrightarrow{\pi} X$ un morphisme plat. Notons $U'\stackrel{{\tilde
j}'}{\hookrightarrow}X'$  l'immersion ouverte $U'=\pi^{-1}(U)\hookrightarrow X'$.  On suppose que
$\pi$ est régulier au dessus de $Z$. Soit $\C$ un champ en groupoïdes sur $U_\et$. Alors, la flèche
de changement de base $$\phi(\C):\pi^*{\tilde j}_*\C\ra {\tilde j}'_*\pi^*\C$$ est une équivalence
dans les deux cas suivants :
\begin{itemize} \item[(i)] $\C$ est discret (c'est-à-dire $\C$
équivalent à un faisceau d'ensembles).  \item[(ii)] $Z$ est $c_2$ et $\C=\tors(U,G)$ avec $G$ un
groupe (ordinaire) fini.
\end{itemize}
\end{theoreme2}

En considérant les fibres, on peut supposer que $\pi$ est un morphisme local de schémas strictement
locaux (la condition $c_2$ ne dépendant que des hensélisés stricts aux points de $Z$).

Soient $x,x'$ les points fermés respectifs de $X,X'$. Par récurrence sur la dimension de $X'$, on
peut supposer que $\phi(\C)_{\bar y'}$ est une équivalence en tout point géométrique $\bar y'$ de
$X'-\{x'\}$ et il suffit de prouver que $\phi(\C)_{x'}$ est une équivalence. On peut de plus
supposer $x\in Z$ (sinon $U=X$ et c'est terminé). Par hypothèse, la fibre spéciale $F=\pi^{-1}(x)$
de $\pi$ est géométriquement régulière.

On a un diagramme commutatif à « carrés » cartésiens (avec des notations un peu abusives)

$$\xymatrix{U'\ar@{^(->}[d]_{j'}\ar@/_3pc/[dd]_{\tilde
j'}\ar[r]&U\ar@{^(->}[d]^{j}\ar@/^3pc/[dd]^{\tilde j}\\
X'-F\ar@{^(->}[d]\ar[r]&X-\{x\}\ar@{^(->}[d]\\ X'\ar[r]^\pi&X}$$

Par hypothèse de récurrence, la flèche de changement de base associée au carré supérieur est une
équivalence de sorte qu'on a une équivalence $\pi^*j_*\C\iso j'_*\pi^*\C$ sur $X'-F$.  Comme $X,X'$
sont strictement henséliens, la flèche de changement de base $\phi(j_*\C)_{x'}$ \begin{eqnarray*}
H^0(X,\tilde j_*\C)&=& H^0(X-\{x\},j_*\C) \\ &\xrightarrow{\pi^*}& H^0(X'-F,\pi^*j_*\C) \\ &=&
H^0(X'-F,j'_*\pi^*\C) \\ &=&H^0(X',\tilde j'_*\pi^*\C) \end{eqnarray*}

s'identifie à la flèche d'image inverse

\begin{equation}\label{20bcphi} \pi^*:\ H^0(X-\{x\},j_*\C) \ra H^0(X'-F,\pi^*j_*\C).  \end{equation}

\begin{lemme2} On peut supposer que $\pi$ est un morphisme essentiellement lisse de schémas
strictement locaux et excellents.  \end{lemme2}

\begin{proof}Notons $\chap{X}$ le complété de $X$ le long de son point fermé et $\widehat{X'}$ le
complété de $X'$ le long de $F$.  Pour tout $S$-espace $\E$ sur $S_\et$ avec $S=X,X'$, on note
$\chap{\E}$ son image inverse sur $\chap{S}$.  On a un diagramme \emph{commutatif}
$$\xymatrix{\widehat{X'}\ar[r]^{\gamma'}\ar[d]^{\hat\pi}&X'\ar[d]^\pi\\\chap{X}\ar[r]^\gamma&X}$$ où
$\gamma,\gamma'$ sont les morphismes de complétion, donc sont plats, et $\hat\pi$ est plat comme
$\pi$ et est un morphisme local de schémas noethériens. Sa fibre spéciale est encore $F$ de sorte
qu'elle est géométriquement régulière. Ainsi, $\hat\pi$ est formellement lisse
(\egalong{IV}{4}{19.7.1}) et donc régulier (\cite{Localisation@Andre}) puisque $\chap{X}$ est local
noethérien complet donc excellent. D'après~\ref{20platc2}, $\chap{Z}=\chap{X}-\chap{U}$ est encore
$c_2$.  D'après le théorème de rigidité de Gabber (\ref{20-prop-rig-champ}) appliqué aux paires
henséliennes $(X,x)$ et $(X',F)$, il suffit, pour prouver que le foncteur~\ref{20bcphi} est une
équivalence, de prouver que le foncteur


\begin{equation}\label{20bcphic} \hat\pi^*:\ H^0(\chap{X}-\{x\},\widehat{j_*\C}) \ra
H^0(\chap{X'}-F,\widehat\pi^*\widehat{j_*\C}) \end{equation} est une équivalence. Mais on a d'une
part $$\widehat{j_*\C}={\chap{j}}_*\chap{\C}$$ d'après la version faisceautique (\ref{20CTC}) du
théorème de rigidité de Gabber,  et, d'autre part $$\chap{\C}=\tors(\chap{U},G)$$ d'après
\cite[III.2.1.5.7]{CNA@Giraud}. Pour prouver que (\ref{20bcphi}) est une équivalence, il suffit donc
de prouver que (\ref{20bcphic}) est une équivalence dans le cas où $X$ est complet, donc excellent
et $\pi$ un morphisme local régulier.






 D'après le théorème de Popescu (\cite{NPD@Swan}), le morphisme régulier $\pi$ est limite projective
filtrante de morphismes locaux essentiellement lisses $\pi_i: X'_i\ra X$. Notons que les $X'_i$ sont
strictement locaux et excellents comme $X$. Comme les $X'_i$ sont cohérents,  le foncteur section
globale commute à la limite projective au sens de \sga{4}{VII}{5.7} de sorte qu'il suffit de prouver
le théorème pour les $\pi_i$.\end{proof}



\subsection{Preuve de~\ref{20rigfaible}} On suppose donc que $\pi$ est un morphisme local
essentiellement lisse de schémas excellents.

Si $\C$ est discret, on invoque le théorème de changement de base par un morphisme (essentiellement)
lisse  pour conclure \sga{4}{XVI}{1.2} : le (\text{i}) du théorème~\ref{20rigfaible} est prouvé.

Supposons donc $\C=\tors(U,G)$. On doit donc prouver pour conclure la preuve  du
théorème~\ref{20rigfaible} la variante suivante du théorème de changement de base lisse de Gabber
(\ref{20-cbl}).

\begin{proposition2}\label{20-cblvar} Considérons un diagramme cartésien
$$\xymatrix{U'\ar@{^(->}[r]^{j'}\ar[d]\ar@{}[rd]|\square&X'\ar[d]^\pi\\ U\ar@{^(->}[r]^{j}&X}$$ où
$\pi$ est un morphisme essentiellement lisse de schémas excellents strictement locaux. On suppose
que le fermé complémentaire $Z=X-U$ est $c_2$ (c'est-à-dire sous ces hypothèses,  que $U$ contient
les points de codimension $1$ (\ref{20remc2})). Alors, le morphisme
$$\pi^*:H^0(X-\{x\},j_*\tors(U,G))\ra H^0(X'-\pi^{-1}\{x\},\pi^*j_*\tors(U,G))$$ est une
équivalence.\end{proposition2}

\begin{proof}Comme $X$ est excellent, le morphisme de normalisation $p:X^\nu\ra X$ est fini. Son
image est donc fermée. Comme $p$ est (ensemblistement) dominant, $p$ est surjectif. Comme $p$ est
surjectif, le foncteur $$j_*\tors(U,G)\ra
j_*p_*p^*\tors(U,G)\stackrel{\hbox{\cite[III.2.1.5.7]{CNA@Giraud}}}{=}p_*j^\nu_*\tors(U^\nu,G)$$ est
fidèle\footnote{On note $\E\mapsto\E^\nu$ le foncteur d'image inverse par $p$.}.  D'après
\ref{20-pertinence-eff} et  le  théorème~\ref{20rigfaible}, (\text{i}), il suffit de prouver que la
flèche

\begin{equation}\label{20imainv}\pi^*:H^0(X-\{x\},p_*j^\nu_*\tors(U^\nu,G))\ra
H^0(X'-\pi^{-1}\{x\},\pi^*p_*j^\nu_*\tors(U^\nu,G))\end{equation} est une équivalence.

Considérons le diagramme cartésien

$$\xymatrix{{X'^\nu}\ar[r]^-{\pi^\nu}\ar[d]^{p'}\ar@{}[rd]|\square&{X}^\nu\ar[d]^p\\
X'\ar[r]^{\pi}&X}$$

Comme $p$ est fini (donc propre), on a $\pi^*p_*=p'_*\pi^{\nu*}$ de sorte que \ref{20imainv}
s'identifie à la flèche d'image inverse

\begin{equation}\label{20imainvbis}\pi^*:H^0(X^\nu-\{x\}^\nu,j^\nu_*\tors(U^\nu,G))\ra
H^0(X'^\nu-(\pi^{\nu})^{-1}\{x\}^\nu,\pi^{\nu*}j^\nu_*\tors(U^\nu,G)).\end{equation}

Notons que, la condition $c_2$ ne dépendant que du normalisé, le complémentaire $Z^\nu$ de $U^\nu$
est encore $c_2$ dans $X^\nu$, et $U^\nu$ contient tous les points de codimension $1$.  On invoque
alors le théorème de changement de base lisse de Gabber \ref{20-cbl}.\end{proof}

\subsection{Comparaison à la complétion : cas des coefficients
abéliens dans le cas non nécessairement nœthérien}

\begin{quote}
Le paragraphe suivant est une \emph{esquisse} de démonstration
de l'analogue du théorème \ref{20rigfaible} pour les
coefficients abéliens. Le cas des schémas nœthériens est traité
dans \cite{TUBULAR@Fujiwara}. Nous reproduisons ici fidèlement
une lettre d'Ofer Gabber aux éditeurs (20 juin 2012).
\end{quote}

\vskip.5cm

Let $(A,I) \to (A',I')$ be a map of henselian pairs with $I$
finitely generated, $I' = I A'$, $\widehat{A} \iso
\widehat{A'}$ ($I$-adic completions). $X = \SP(A)$, $X' =
\SP(A')$, $\pi : X' \ra X$, $U = X - V(I)$, $U' = X' -
V(I')$, $j : U \to X$, $j' : U' \ra X'$.

\begin{quote}
\textbf{CTC} : For every torsion abelian  sheaf $F$ on $U$,
the base change arrow $\pi^*R^qj_*F \to R^qj'_*\pi^*F$ is an
isomorphism for all $q$.
\end{quote}

\textit{Analogue of \ref{20rigfaible}} (notations as there) : If $F$
is a sheaf of $\ZZ/n\ZZ$-modules on $U$ where
$n >0$ is invertible on $X$, then $\pi^*R^qj_*F \ra R^qj'_*
\pi^*F$ are isomorphisms.

This is reduced to CTC by the same argument.

\textit{Sketch of proof of CTC using Zariski-Riemann spaces}
: For comparing stalks we may assume $A$, $A'$ strictly
henselian and $I$ a proper ideal, and we want $$ H^q(U,F)
\iso H^q(U',F). \leqno (*) $$ We call a finitely generated
ideal $J \subset A$ containing a power of $I$
\textit{admissible}. We consider the admissible blow-ups
$\Bl_{J}(X)$ which form a cofiltered category using
$X$-scheme morphisms. In general there can be more than one
$X$-morphism between two admissible blow-ups but if we
restrict ourselves to $J$'s with $V(J) = V(I)$ (set
theoretically) (so that $U$ is schematically dense in the
blow-up), there is at most one. Define $J \le J'$ iff there
is an $X$-morphism $\Bl_{J'}(X) \ra
\Bl_J(X)$. This is a filtered preorder. When $V(J)
\subset V(J')$, $J \le J'$ is equivalent to the condition
that for some $n >0$ and ideal $K$, $J'^n = JK$. Thus we
have an isomorphism of the preordered set of admissible
$J$'s of full support in $A$ and the corresponding set for
$A'$. Let $ZRS_I(X) = \varprojlim \Bl_J(X)$ (a
locally ringed space). For the closed point $s$ of $X$ we
can consider the special fiber $ZRS_I(X)_s$ and its étale
topos, which for our purposes may be defined as the
projective limit of the étale topoi
$(\Bl_J(X)_{s})_{\et}$ as in
[\textbf{SGA 4}]. It has enough points by Deligne's theorem.
The points are given by "geometric points" of $ZRS_I(X)_s$
(i. e. a point and a choice of a separable closure of the
residue field). For every admissible $J$ we have $$ j_J : U
\hookrightarrow \Bl_J(X) $$ giving a spectral
sequence (using proper base change) $$
H^p(\Bl_J(X)_s, R^qj_{J*}F) \Rightarrow
H^{p+q}(U,F). \leqno (**) $$ We pass to the limit using the
general theory of \sgacourt{4}{vi}. We get a spectral sequence
$(**)_{\mathrm{lim}}$ involving cohomology on
$(ZRS_J(X)_s)_{\et}$. Since the latter topos is the
same for $X'$, to show $(*)$ we use the morphism of the limit
spectral sequence to reduce to stalks of the limits of the
$R^qj_{J*}$ sheaves.

Using the study in \cite{TUBULAR@Fujiwara} of the local rings of ZRS's
and their henselizations, one reduces $(*)$ to the case of
local rings at geometric points of the special fibers of
ZRS's.  Thus we are reduced to the case $A$, $A'$ are
henselian and $I$-valuative. Say $I = (\varphi)$. Then
$A[\varphi^{-1}]$ is a henselian local ring with maximal
ideal corresponding to $P = ⋂ I^n$, and $A/P$ is a
henselian valuation ring whose valuation topology is the
$\varphi$-adic one. In this case to prove $(*)$ one reduces to
the corresponding statement for $\mathrm{Frac}(A/P) \to
\mathrm{Frac}(A'/P')$. In fact for $K \to K'$ a dense
embedding of henselian valued fields, if we choose separable
closures $K_{\mathrm{sep}}$, $K'_{\mathrm{sep}}$ and a map
between them we have $\mathrm{Gal}(K'_{\mathrm{sep}}/K')
\iso \mathrm{Gal}(K_{\mathrm{sep}}/K)$, using forms of
Krasner's lemma (cf. \accourt{VI}{8, exercices 12, 14 a}).

Note : For admissible $J$, $\Bl_J(X') \ra
\Bl_J(X)$ gives an isomorphism on $I$-adic
completions (as in the discussion in the proof of
\ref{20-commut-princ}) as for every $m$ the map $$ \oplus_n
J^n \to \oplus_n J^nA' $$ is an isomorphism mod $I^m$.



\section{Rigidité de la ramification II : forme forte}

\subsection{Générisations étales immédiates} On note $X_{(x)},X^\h_{(x)}$ et $\chap{X}_{(x)}$ les
localisés, hensélisés et complétés respectivement de $X$ en $x$. On note $\overline{\{y\}}$
l'adhérence de $y$ dans $X$ munie de sa structure réduite. L'hensélisation et la complétion
commutent aux immersions fermées de sorte que $\overline{\{y\}}^\h,\chap{\overline{\{y\}}}$
respectivement coïncident avec l'image inverse de $\overline{\{y\}}$ par les morphismes
d'hensélisation, complétion respectivement.  Rappelons
(\refext{XIV}{14-definition-specialisation-immediate}) qu'une générisation $y\in X$ d'un point $x$
d'un schéma $X$ est une \textit{générisation étale immédiate} de $x$ si l'hensélisé strict en $\bar
x$ de l'adhérence de $y$ a une composante irréductible de dimension $1$.






\begin{lemme2}\label{20et-imm} Soit $y$ une générisation de $x$. Notons $c:\chap{X}_{(x)}\ra
X_{(x)}$ le morphisme de complétion. Alors, $y$ est une générisation étale immédiate de $x$ si et
seulement si l'un des points maximaux de $c^{-1}(y)$ est une générisation étale immédiate du point
fermé de $\chap{X}_{(x)}$.\end{lemme2}

 \begin{proof} Notons pour simplifier $Y=\overline{\{y\}}$. Observons d'abord qu'un des trois
schémas $Y_{(x)},Y^\h_{(x)}$ et $\chap{Y}_{(x)}$ possède un point maximal de dimension nulle si et
seulement si chacun est réduit (ensemblistement) à son point fermé. On peut donc exclure ce cas. Le
morphisme $Y_{(\bar x)}\ra Y^\h_{(x)}$ est fidèlement plat et entier. Donc,  l'hensélisé strict
possède un point maximal de dimension $1$\footnote{On devrait plutôt dire point maximal dont
l'adhérence est de dimension $1$.} si et seulement si l'hensélisé $Y^\h_{(x)}$ possède un point
maximal de dimension $1$. D'après (\ref{20propc2}), $Y^\h_{(x)}$ possède un point maximal de
dimension $1$ si et seulement si $\chap{Y}_{(x)}(\ref{20propc2})$ possède un point maximal de
dimension $1$. Par platitude de $c$, il s'envoie nécessairement sur $y$, le point générique de
$Y$.\end{proof}

On peut caractériser agréablement les générisations étales immédiates.


\begin{lemme2} Soit $f:X_{(\bar x)}\ra X_{(x)}$ le morphisme d'hensélisation strict. Les
générisations étales immédiates de $x$ sont les images $y=f(y')$ des $y'\in X_{(\bar x)}$ tels que
$\dim\overline{\{y'\}}=1$.\end{lemme2}




\begin{proof} Soit $y'\in X_{(\bar x)}$ tel que $\dim\overline{\{y'\}}=1$. L'image $y=f(y')$ est une
générisation stricte de $x$ (car par exemple les fibres de $f$ sont discrètes). Pour cette même
raison, $\overline{\{y'\}}$ est une composante de $f^{-1}(\overline{\{y\}})=\overline{\{y\}}_{(\bar
x)}$. Inversement, si $y$ est une générisation étale immédiate de $x$, le point générique $y'$ d'une
composante de dimension $1$ de $\overline{\{y\}}_{(\bar x)}$ s'envoie sur $y$ (platitude de $f$) et
son adhérence est de dimension $1$.\end{proof}




\begin{exemple2}Prenons l'exemple du pincement de \egalong{IV}{2}{5.6.11}. En conservant les
notations de \textit{loc. cit.}, l'anneau pincé $C$ est local noethérien de dimension $2$ et son
normalisé a deux idéaux maximaux de hauteur $1,2$ respectivement. D'après~\ref{20hnu=nuh},
l'hensélisé de $C$ a deux composantes irréductibles de dimension $1$ et $2$ de points génériques
$c,c'$. Comme dans la preuve de \refext{XIV}{14-prop-Laszlo-2}, , ceci assure l'existence de $\bar
c$ (au dessus de $c$) dans l'hensélisé strict de $C$ dont l'adhérence est de dimension $1$ et donc
que le point générique de $\Spec(C)$ est une générisation étale immédiate de son point
fermé.\end{exemple2}


\subsection{Couples associés et condition (*)}\label{20condstar} Commençons par une définition.

\begin{definition2}\label{20defass}Soit $x$ un point d'un schéma $X$. Choisissons une clôture
séparable de $k(x)$ définissant un point géométrique $\bar x$ de $X$.  \begin{itemize} \item[(i)]
Soit $G$ un schéma en groupes sur $X$. On définit les sections locales de $G$ à support dans $\bar
x$ par la formule $$H^0_{\bar x}(G)=\Ker(H^0(X_{(\bar x)},G)\ra H^0(X_{(\bar x)}-\{\bar x\},G)).$$
\item[(ii)] Soit $\C$ un champ (en groupoïdes)  sur $X$ et $p$ un nombre premier. On dit que $(x,p)$
est associé de $\C$ et on écrit $(x,p)\in \Ass(\C)$ si il existe $\sigma\in\C_{\bar x}$ tel que
$H^0_{\bar x}(\aut(\sigma))$ ait de la $p$-torsion.  \item[(iii)] Soit $\C$ un champ ind-fini (en
groupoïdes) sur un ouvert $U$ de $X$. On dit que $\C$ vérifie la condition (*) si pour tout $x\in
X-U$ de caractéristique $p>0$, il n'existe pas de générisation étale immédiate $y$ de $x$ telle que
$(p,y)\in\Ass(\C)$.  \end{itemize} \end{definition2}

Remarquons que la condition $(x,p)$ associé ne dépend pas du choix de la clôture séparable de
$k(x)$.

\begin{exemple2}\label{20exetoile} Supposons $X$ normal et $G$ groupe fini. Soit $U$ un ouvert de
$X$. Alors, $(x,p)$ est associé de $\C=\tors(U,G)$ si et seulement si $p|\card(G)$ et $x$ est un
point maximal de $U$. En effet, l'unique objet de $\C_{\bar x}$ est le torseur trivial $\sigma$ et
$\aut(\sigma)=G$. Or, $X_{(\bar x)}-\{\bar x\}$ est connexe (resp. vide) si $x$ non maximal (resp.
maximal). Ainsi, on a $H^0_{\bar x}(\aut(\sigma))=\{1\}$ (resp. $H^0_{\bar x}(\aut(\sigma))=G$). On
déduit que $\C$ vérifie (*) si et seulement si $U$ contient tous les points de codimension $1$ dont
la caractéristique divise l'ordre de $G$.  \end{exemple2}

\begin{lemme2}\label{20ass-comp} Soit $f:X\ra Y$ un morphisme plat de schémas noethériens, $x\in X$
d'image $y=f(x)$ dans $Y$  et $\C$ un champ en groupoïdes sur $Y$. Alors, $(x,p)\in\Ass(f^*\C)$ si
et seulement si $(y,p)\in\Ass(\C)$ et $x\in\Max(f^{-1}(y))$.  \end{lemme2}

\begin{proof} Choisissons un point géométrique $\bar x$ au dessus de $x$, qui définit $\bar y$ au
dessus de $y$.

Supposons $(x,p)\in\Ass(f^*\C)$.  Comme la flèche $(f^*\C)_{\bar x}\ra\C_{\bar y}$ est une
équivalence, il existe $\sigma\in \C_{\bar y}$ et $g\in\Aut(\sigma)$ tel que $f^*g$ est
d'ordre de $p$ et de  support $\{\bar x\}$. Notons $F$ l'hensélisé strict de $f^{-1}(y)$ en $\bar
x$. C'est aussi la fibre de l'hensélisé $\varphi: X_{(\bar x)}\ra Y_{(\bar y)}$ de $f$ au dessus de
$\bar y$. Si $F$ n'était pas réduit à $\bar x$, un des points de $F$ ne serait pas dans le support
de $f^*g$ de sorte que $f^*g$ serait l'identité en ce point. Mais $f^*g$ est constant sur
$F=\varphi^{-1}(\bar y)$ de sorte que $f^*g$ serait l'identité également en $\bar x\in F$, ce qui
n'est pas. Donc, $F$ est réduit à $\bar x$ de sorte que $\dim \OO_{f^{-1}(y),x}=0$ (puisqu'un anneau
local a même dimension que son hensélisé strict) et $x\in\Max(f^{-1}(y))$. De plus, $g$ est trivial
sur $\varphi(X_{(\bar x)}-\{\bar x\})=Y_{(\bar y)}-\{\bar y\}$ (fidèle platitude de $\varphi$) ce
qui assure $(y,p)\in\Ass(\C)$.


Inversement, supposons  $(y,p)\in\Ass(\C)$ et $x\in\Max(f^{-1}(y))$. On a donc un automorphisme $g$
d'ordre $p$ de $\sigma\in\C_{\bar y}$ de support $\{\bar y\}$. Le support de $\varphi^*g$ est la
fibre $\phi^{-1}(\bar y)=f^{-1}(y)_{(\bar x)}$. Comme $x$ est maximal dans $f^{-1}(y)$, on déduit
(dimension) que le schéma local $\phi^{-1}(\bar y)$ est de dimension nulle donc réduit à $\bar x$,
ce qu'on voulait.  \end{proof}


\begin{corollaire2}\label{20corass-comp} Soit $(X,x)$ un schéma local noethérien hensélien,
$U=X-\{x\}$ l'ouvert complémentaire du point fermé et $c:\chap{X}\ra X$ le morphisme de complétion.
Alors, le champ en groupoïdes $\C$ sur $U$ vérifie (*) si et seulement $\chap{\C}=c^*\C$ vérifie
(*).\end{corollaire2}

\begin{proof} On note encore $x$ le point fermé de $\chap{X}$ et on choisit un point géométrique
$\bar x$ au dessus de $x$.

Supposons que $\chap{\C}$ vérifie (*). Soit $(y,p=\car(x))\in\Ass(\C)$ et notons $Y$ l'adhérence de
$y$ dans le localisé (Zariski) de $X$ en $x$. Il s'agit de montrer que toutes les composantes de
$Y_{(\bar x)}$ sont de dimension $\geq 2$, ou encore (\ref{20propc2}) que toutes les composantes de
$\chap{Y}=\chap{Y}_{(x)}^h$ sont dimension $\geq 2$, c'est-à-dire (\ref{20propc2} à nouveau), que
toutes les composantes de $\chap{Y}_{(x)}$ sont de dimension $\geq 2$. Mais c'est bien le cas car,
d'après le lemme~\ref{20ass-comp}, on a $(\hat y,p)\in\Ass(\chap{\C})$.

Inversement, supposons que $\C$ vérifie (*). Soit donc $(\hat y,p=\car(x))\in\Ass(\chap{\C})$ et
soit $y=c(\hat y)$ et supposons que l'adhérence de $\hat y$ est de dimension $1$.  Si une des
composantes  $Y=\overline{\{y\}}$ était de dimension $>1$, toutes les composantes de
$\chap{Y}=c^{-1}(Y)$ seraient de dimension $>1$ (\ref{20propc2}). D'après le lemme~\ref{20ass-comp},
on sait que $\hat y$ est maximal dans $c^{-1}(y)$ donc dans $\chap{Y}$ de sorte que
$\dim\overline{\{\hat y\}}>1$, une contradiction. \end{proof}

\subsection{Le théorème de rigidité de la ramification}

\begin{theoreme2}[Rigidité de la ramification II]\label{20rigfort} Soit $\pi:X'\ra X$ un morphisme
plat de schémas noethériens, régulier au dessus d'un sous-schéma fermé $Z\subset X$.  Soit
$j:U=X-Z\hookrightarrow X$ l'immersion ouverte du complémentaire de $Z$ et $\C$ un champ ind-fini
sur $U$ vérifiant la condition (*). Alors, la flèche de changement de base $\pi^*j_*\C\ra
j'_*\pi'^*\C$ est une équivalence.\end{theoreme2}

\begin{proof} D'après le théorème de rigidité de la ramification I (\ref{20rigfaible}), le théorème
est vrai dans le cas discret de sorte que $\pi^*j_*\C\ra j'_*\pi'^*\C$ est toujours pleinement
fidèle.  Comme dans la preuve de \ref{20rigfaible}, on peut supposer $X,X'$ strictement locaux de
point fermés $x,x'$ et $\pi$ morphisme local. Par récurrence sur la dimension de $X$, on peut
supposer que le changement de base par $\pi$ est une équivalence pour l'immersion $U\hookrightarrow
X-\{x\}$ de sorte qu'on peut supposer $U=X-\{x\}$. Comme dans la preuve de \ref{20rigfaible} et en
utilisant l'invariance par complétion de la condition (*) (\ref{20corass-comp}), on peut supposer de
plus $X$ complet et $\pi$ morphisme essentiellement lisse et local et il s'agit de démontrer que la
flèche $$\pi^*:H^0(X-\{x\},\C)\ra H^0(X'-\pi^{-1}\{x\},\C')$$ est essentiellement surjective
(puisque en tout point $y'\not=x'$ la fibre du changement de base est une équivalence par
hypothèse).

Soit donc $\sigma'$ un objet de $H^0(X'-\pi^{-1}\{x\},\C')$. La condition (*) étant stable par
passage aux sous-gerbes (maximales), on peut comme dans la preuve de \ref{20-red-gerb} en
considérant la sous-gerbe maximale de $\C'$ engendrée par $\sigma'$,  supposer de plus que $\C$ est
une gerbe. Comme $\C$ est ind-finie, on peut supposer que $\C$ est \emph{constructible}.

Pour tout point maximal $y\in U$, notons $i_y$ le morphisme canonique $$i_y:\Spec(k(y))\ra
U=X-\{x\}.$$


 Soit $$\Psi : \C\ra\D:=\prod_{y\in \Max(U)}i_{y*}i_{y}^*\C$$ le morphisme déduit des morphismes
d'adjonction. La catégorie fibre de $i_{y*}i_{y}^*\C$ sur un ouvert étale $V\ra X$ s'identifie aux
sections rationnelles de $\C$ définies au voisinage (Zariski) des points maximaux de $V$ au dessus
de $y$. On déduit que $\Psi$ est conservatif et couvrant (étale localement surjectif sur les flèches
et les objets). Pour toute section $\tau\in H^0(U,\D)$ (vu comme un morphisme de $U$-espaces $\tau
:U\ra\D$), le champ des relèvements $K(\tau)=U\times_{\D}\C$ associé  est  donc une gerbe
(\cite[IV.2.5.4]{CNA@Giraud}), évidemment constructible.

Il suffit alors (exercice) de vérifier que la flèche de changement de base est essentiellement
surjective pour \begin{itemize}

  \item [1)] les gerbes $\G=i_{y*}i_{y}^*\C$; \item [2)] la gerbe des relèvements
$K(\tau)=U\times_{\D}\C$ associée à $\tau\in H^0(U,\D)$.  \end{itemize}

\subsubsection{Premier cas : changement de base pour
$\G=i_{y*}i_{y}^*\C$}

Supposons donc $\G=i_{y*}i_{y}^*\C$. Quitte à changer $X,U$ en $\overline{\{y\}},
U\cap\overline{\{y\}}$, on peut supposer que $X$ est irréductible de point générique $y$.

Si la dimension de $X$ est  $1$, on a $U=\{y\}$ et $\G_{\bar y}=\tors(\Spec(\overline{k(y)}),G)$
avec  $p=\car(y)$ ne divisant pas l'ordre de $G$ (cf. l'argument dans l'exemple~\ref{20exetoile}).
On invoque alors le changement de base par un morphisme lisse usuel (\cite[VII.2.1.2]{CNA@Giraud}).

On suppose donc que la dimension de $X$ est $>1$.

Choisissons une clôture séparable $k(y)\hookrightarrow k_y$ et notons ${j_y}:\Spec(k_y)\ra U$ est le
morphisme canonique. On a $$j_y^*\C=\tors(\Spec(k_y),G)$$ où $G$ est un groupe fini constant. Comme
$$i_{y*}i_y^*\C\ra {j_y}_*j_y^*\C$$ est fidèle, on peut (\ref{20-pertinence-eff}) remplacer $\G$ par
${j_y}_*j_y^*\C={j_y}_*\tors(\Spec(k_y),G)$.

\begin{lemme2} On a $R^1j_{y*}G=\{*\}$ et
${j_y}_*\tors(\Spec(k_y),G)=\tors(U,{j_y}_*G)$.\end{lemme2} \begin{proof} La seconde égalité découle
de la première et de la formule (\cite[V.3.1.5]{CNA@Giraud})
$$\pi_0({j_y}_*\tors(\Spec(k_y),G))=R^1j_{y*}G.$$ Soit $\tilde A$ l'hensélisé strict de $X=\Spec(A)$
en un point géométrique $\bar\xi$ de $X$. C'est une limite inductive filtrante d'algèbres $A_i$ de
type finies qui sont génériquement étales. On déduit que $j_y^{-1}(\Spec(\tilde A))$ est le spectre
de la limite inductive filtrante des algèbres étales $B_i=k_y\otimes_{k(y)}A_i$ qui sont donc
scindées puisque $k_y$ est séparablement clos. Ainsi, les schémas considérés étant cohérents, on a
$$(R^1j_{y*}G)_{\bar\xi}=H^1(j_y^{-1}(\Spec(\tilde A)),G)=\varinjlim
H^1(\Spec(B_i),G)=\{*\}.$$\end{proof}

En terme de module galoisien, $j_{y*}G$ est l'induite (continue) $\Hom_c(\Gamma,G)$ où
$\Gamma$ est le groupe profini $\Gal(k_y/k(y))$. En écrivant $\Gamma=\limind
\Gal(K_\alpha/k(y))$ où $(K_\alpha/k(y))_\alpha$ est le système inductif des sous-extensions
galoisiennes finies de $k_y/k(y)$, on trouve $$j_{y*}G=\limind j_{\alpha*}G$$ où
$j_\alpha:\Spec(K_\alpha)\ra\Spec(k(y))\ra U$ est le morphisme canonique. Comme $U, U'$ sont
noethériens donc cohérents, on a $$H^0(U,\tors(U,\limind j_{\alpha*}G))=\limind H^0(U,\tors(U,
j_{\alpha*}G))$$ et $$H^0(U',\pi^*\tors(U,\limind j_{\alpha*}G))=H^0(U',\tors(U',\pi^*\limind
j_{\alpha*}G))=\limind H^0(U,\tors(U', \pi^*j_{\alpha*}G))$$ de sorte qu'on est réduit pour le cas
1) à étudier le changement de base pour la gerbe $\G_\alpha=\tors(U,j_{\alpha*}G)$.


Soit $p:W\ra X$ la normalisation de $X$ dans $\Spec(K_\alpha)\ra X$ : c'est un morphisme fini (car
$X$ est excellent) et surjectif de sorte que  $W$ est semi-local et hensélien (comme $X$). On déduit
que $W$ est la réunion disjointe de ses hensélisés aux points fermés. Comme $W$ est intègre,  $W$
est  strictement local : on note $w$ son point fermé $w$. De plus, $W$ est normal, donc
géométriquement unibranche de sorte que $j_{\alpha *}G=p_*G_{|W-w}$. Comme $R^1p_*G$ est trivial
($p$ est fini), on déduit l'égalité $$\G_\alpha=p_*\tors(W-\{w\},G)$$ comme dans la preuve du
lemme~\ref{20pdirtors} \textit{infra}.


 En utilisant le changement de base propre pour $p$, on est ramené à prouver que la flèche
$$\Tors(W-\{w\},G)\ra\Tors(W'-\pi^{-1}\{w\},G)$$ est une équivalence. Deux cas se
présentent.

Rappelons qu'on a supposé que la dimension de $X$ (ou $W$, c'est la même chose) est  $>1$. Dans ce
cas, $\{w\}$ est $c_2$ dans $W$ et on invoque la variante du théorème de changement de base lisse de
Gabber (\ref{20-cblvar}).

\subsubsection{Deuxième cas : changement de base pour la gerbe des
relèvements $K=K(\tau)$}

Comme le morphisme $K(\tau)\ra \C$ est visiblement fidèle (\cite[IV.2.5.2]{CNA@Giraud}), la gerbe
des relèvements vérifie (*) comme $\C$. Comme $\Psi_{\bar y}$ est une équivalence pour tout point
maximal $y\in U$, on déduit que $K_{\bar y}$ est la gerbe triviale (autrement dit équivalente au
faisceau d'ensembles ponctuels) en tous ces points. Par hypothèse de récurrence, il suffit pour
achever la preuve de prouver le lemme suivant.

\begin{lemme2} Il existe une immersion fermée
$i:F\subset X$ nulle part dense telle que $K=i_*i^*K$.
\end{lemme2}

\begin{proof}Il suffit de prouver que pour tout $y$ maximal, il existe un ouvert de Zariski
contenant $y$ sur lequel $K$ est triviale. Par construction, il existe un voisinage étale $V\ra X$
de $y$ et $\sigma\in K(V)$. Comme $\underline{\Aut}(\sigma)$ est un faisceau constructible de
$V_\et$, l'isomorphisme $$\{\Id\}_{\bar y}\iso\underline{\Aut}(\sigma)_{\bar y}$$ provient
d'un isomorphisme $$\{\Id\}_{W}\iso\underline{\Aut}(\sigma)_W$$ sur un voisinage étale $W\ra
V\ra X$ de $y$. Quitte à localiser, on peut supposer que $W\ra X$ est un revêtement galoisien de son
image $U\subset X$. La section $\sigma$ descend sur $U$ et n'a pas d'automorphisme par construction,
ce qu'on voulait.  \end{proof} \end{proof}



\section{Appendice 1 : sorites champêtres} D'après \ref{20-prop-rig} i), on sait déjà que le
foncteur $H^0(U,\C)\ra H^0(\chap{U},p^*\C)$ est pleinement fidèle que $\C$ soit ind-fini ou non.

Soit $\C$ un champ ind-fini sur $Y$. On cherche des conditions assurant que l'hypothèse

\begin{hypothese}\label{20-hypo-champ}Soit $f:X\ra Y$ un morphisme de schémas. On suppose que pour
tout faisceau d'ensembles $\F$ sur $Y$, la flèche $H^0(Y,\F)\ra H^0(X,f^*\F)$ est
bijective.\end{hypothese}

entraîne que la conclusion

\begin{conclusion}\label{20-concl-champ}  La flèche $\phi: H^0(Y,\C)\ra H^0(X,f^*\C)$ est une
équivalence de catégories.\end{conclusion}

est vraie, autrement dit assurant que l'assertion

\begin{assertion}\label{20-cond-champ}On a l'implication
$\ref{20-hypo-champ}\Rightarrow\ref{20-concl-champ}$.  \end{assertion}

est vraie. On sait déjà que \ref{20-hypo-champ} entraîne que $\phi$ est pleinement fidèle (cf.
\ref{20-sec-enonc}).

\subsection{Premières réductions}

Commençons par un lemme formel :

\begin{lemme2}\label{20-pertinence-eff} Soit $f:X\ra Y$ un morphisme de schémas et $\C_1\ra\C_2$ un
morphisme de champs sur $Y$ qu'on suppose fidèle. Si \ref{20-cond-champ} est vraie pour $\C_2$,
alors \ref{20-cond-champ} est vraie pour $\C_1$.\end{lemme2}

\begin{proof} On a déjà observé (\ref{20-sec-enonc}) que $\phi$ est pleinement fidèle. Soit donc
$$c_1^X\in H^0(X,f^*\C_1)=\Hom_X(X,f^*\C_1)$$ dont on cherche un antécédent dans $H^0(Y,\C_1)$. Son
image $$c_2^X\in H^0(X,f^*\C_2)$$ a un antécédent (à isomorphisme près) $$c_2^Y\in H^0(Y,\C_2).$$

Le couple $(c_1^X,c_2^X=f^*c_2^Y)$ définit une section de la gerbe des relèvements
$$K(f^*c_2^X)=X\times_{f^*\C_1}f^*\C_2.$$

D'après \cite[V.1.4.3]{CNA@Giraud}, la gerbe $K(f^*c_2^X)=X\times_{f^*\C_1}f^*\C_2$ s'identifie à
$$f^*K(c_2^X)=f^*(Y\times_{\C_1}\C_2).$$ Or, $K(c_2^X)$ est un faisceau d'ensembles car $\C_1\ra
\C_2$ est fidèle. Donc, $(c_1^X,c_2^X=f^*c_2^Y)\in H^0(X,f^*K(c_2^X))$ a un unique antécédent de la
forme $(c_1^Y,c_2^Y)$ et $c_1^Y$ est bien l'antécédent cherché. \end{proof}

\begin{lemme2}\label{20-red-gerb} Si \ref{20-cond-champ} est vrai pour toute gerbe (resp. toute
gerbe ind-finie), alors \ref{20-cond-champ} est vrai pour tout champ (resp. tout champ
ind-fini).\end{lemme2}

\begin{proof} Soit $t\in H^0(X,f^*\C)$ et $\gamma_t\subset f^*\C$ la sous-gerbe maximale engendrée
par $t$ dans $f^*\C$ (\cite{CNA@Giraud}, III.2.1.3.2)\footnote{Dans \textit{loc. cit.}, $\pi_0(\C)$
est noté $\mathrm{Ger}(\C)$, qui n'est pas actuellement la notation standard.}.  Ceci définit une
section $\tau\in\pi_0(f^*\C)$ du faisceau d'ensembles $\pi_0(f^*\C)$ des sous-gerbes maximales de
$f^*\C$ (\textit{loc. cit.}, 2.1.4). D'après \textit{loc. cit.}, 2.1.5, la flèche naturelle
$$\pi_0(f^*\C)\ra f^*\pi_0(\C)$$ est bijective. Mais, par hypothèse, la flèche $$H^0(Y,\pi_0(\C))\ra
H^0(X,f^*\pi_0(\C))=H^0(X,\pi_0(f^*\C))$$ est bijective de sorte qu'il existe une (unique)
sous-gerbe (maximale) $\gamma\subset\C$ telle que $f^*\gamma=\gamma_t$, qui sera ind-finie si $\C$
l'est.  L'image dans $H^0(Y,\C)$ de l'antécédent de $t\in H^0(X,f^*\gamma)$ dans $H^0(Y,\gamma)$ est
l'antécédent cherché.\end{proof}

\subsection{Réduction au cas d'un champ de torseurs sous un groupe fini constant} Admettons pour un
instant le résultat suivant, généralisation au cas des champs de la résolution flasque de Godement.

\begin{lemme2}[Lemme d'effacement]\label{20-lemme-efface} Soit $\gamma$ une gerbe ind-finie sur un
schéma cohérent $X$. Il existe un groupe ind-fini $\G$ sur $X$ et un foncteur fidèle
$\gamma\hookrightarrow \underline{\Tors}(X,\G)$.\end{lemme2}

On peut alors prouver le critère suivant

\begin{proposition2}\label{20-prop-red-champ} Soit $f:X\ra Y$ un morphisme de schémas cohérents.  On
suppose que pour tout faisceau d'ensembles $\F$ sur $Y$, la flèche $H^0(Y,\F)\ra H^0(X,f^*\F)$ est
bijective (\ref{20-hypo-champ}).  On suppose en outre  que pour tout morphisme fini $p:Y'\ra Y$
induisant $f':X'=X\times_Y Y'\ra Y'$ et tout groupe fini constant $G$, la flèche
$\Tors(Y',G)\ra\Tors(X',G)$ est une équivalence.  Alors, pour tout champ ind-fini $\C$ sur $Y$, la
flèche $H^0(Y,\C)\ra H^0(X,f^*\C)$ est une équivalence.\end{proposition2}

\begin{proof} Seule l'essentielle surjectivité pose problème.  Les lemmes d'effacement,
\ref{20-pertinence-eff} et \ref{20-red-gerb} permettent de supposer que $\C=\underline{\Tors}(Y,\G)$
où $\G$  est un groupe ind-fini sur $Y$. Comme $X,Y$ sont cohérents, la cohomologie non abélienne
commute aux limites inductives filtrantes \sga{4}{VII}{remarque VII.5.14}. Comme $\G$ est ind-fini,
il est limite inductive filtrante de faisceaux en groupes constructibles \sga{4}{IX}{2.7.2} : on
peut donc supposer $\G$ constructible. Puisque $Y$ est cohérent,  il existe (\textit{loc. cit.},
2.14) une famille finie de morphismes finis $p_i:Y_i\ra Y$ et des groupes finis constants $G_i$ tels
que $\G$ se plonge dans le produit $\prod p_{i*}G_i$. On a donc un morphisme fidèle
$$\underline{\Tors}(Y,\G)\hookrightarrow \underline{\Tors}(Y,\prod{p_{i*}}G_i)=
\prod\underline{\Tors}(Y,{p_{i*}}G_i)$$ grâce à \cite{CNA@Giraud}, III.2.4.4.

Utilisant à nouveau \ref{20-pertinence-eff}, on peut supposer $$\C=\underline{\Tors}(Y,{p_*}G)$$
avec $G$ groupe fini constant et $p:Y'\ra Y$ fini.


\begin{lemme2}\label{20pdirtors}On a
$\underline{\Tors}(Y,{p_*}G)=p_*\underline{\Tors}(Y',G)$.\end{lemme2}

\begin{proof}Comme $p$ est fini, $R^1p_*G$ est trivial. Mais
$\pi_0(p_*\underline{\Tors}(Y',G))=R^1p_*G$ (\cite{CNA@Giraud}, V.3.1.9.1) de sorte que
$p_*\underline{\Tors}(Y',G)$ est une gerbe, visiblement neutre et vaut donc nécessairement
$\underline{\Tors}(Y,{p_*}G)$.\end{proof}

Le théorème de changement de base propre pour les faisceaux (trivial dans ce cas) assure  qu'on a
$f^*p_*G=p'_*f'^*G=p'_*G$. La flèche $$H^0(Y,p_*\underline{\Tors}(Y',G))\ra
H^0(X',f^*p_*\underline{\Tors}(Y',G))$$ s'identifie alors à la flèche naturelle


$$\begin{array}{llll} \Tors(Y',G)&=& H^0(Y,p_*\underline{\Tors}(Y',G))\\
&\ra&H^0(X,f^*p_*\underline{\Tors}(Y',G)) \\ &{=}& H^0(X,\underline{\Tors}(X,f^*p_*G))&\
(\text{d'après \ref{20pdirtors} et \cite[III.2.1.5.7]{CNA@Giraud}}) \\ &=&
H^0(X,\underline{\Tors}(X,p'_*G)) \\ &=& H^0(X',p'_*\underline{\Tors}(X',G)) \\ &{=}&
H^0(X',\underline{\Tors}(X',G)) \\ &=& \Tors(X',G) \end{array}$$ qui est bijective par
hypothèse.\end{proof}



\subsection{Preuve du lemme d'effacement} Soit $X$ un schéma cohérent.

\begin{lemme2}\label{cres-alg-clos}Il existe un schéma affine $X'$, un morphisme quasi-compact et
surjectif $f:X'\ra X$ tel que pour tout $x'\in X'$, le corps résiduel $k(x')$  est la clôture
algébrique du corps résiduel $k(f(x))$. \end{lemme2}

\begin{proof} Comme $X$ est quasi-compact, on peut recouvrir $X$ par un nombre fini d'ouverts
affines $X_i$. Le morphisme $f:\sqcup X_i\ra X$ est surjectif et quasi-compact ($X$ est
quasi-séparé). Comme toutes les extensions résiduelles sont des isomorphismes, on peut donc supposer
$X=\Spec(A)$ affine quitte à changer $X$ en  $\sqcup X_i$.




\begin{sslemme2}\label{20-c-co} Soient $I$ l'ensemble des polynômes unitaires (non constants) de
$A[X]$ et $$T'(A)=A[X_P, P\in I]/(P(X_P))$$  et $$f:X'=\Spec {T'(A)}\ra X=\Spec A.$$ Le morphisme
$f$ est surjectif et, pour tout $\xi\in X'$, le corps résiduel $k(\xi)$ est la clôture algébrique de
$k(f(\xi))$.\end{sslemme2}

\begin{proof} Soit $x\in X$ et $\overline{k(x)}$ une clôture algébrique de $k(x)$. La fibre
schématique $f^{-1}(x)$ est le spectre de $$B=k(x)[X_P, P\in I]/(\tilde P(X_P))$$ où $\tilde P$
désigne l'image de $P\in I$ par le morphisme de localisation des coefficients $$A[X]\rightarrow
k(x)[X].$$ Le choix de racines $x_P\in \overline{k(x)}$ pour tous les polynômes $P$ de $I$ définit
un point de $f^{-1}(x)$ assurant la surjectivité de $f$.

Soit alors $\xi\in f^{-1}(x)$ : c'est un point fermé car $f$ est entier et $k(\xi)$ est algébrique
sur $k(x)$. Soit $Q$ un polynôme unitaire de $k(x)$ de degré $d>0$. Il existe $a\in A$ d'image
$a(x)$ non nulle dans $k(x)$, un polynôme unitaire $P\in I$ et un entier $n>0$ tel que
$$s^{nd}Q(X)=P(s^nX).$$ On déduit que l'image $X_P/s^n$ dans $k(\xi)$ est une racine de $Q$. Comme
$k(\xi)$ est algébrique sur $k(x)$, ceci assure que $k(\xi)$ est une clôture algébrique de $k(x)$.
(C'est un exercice (facile) de théorie de Galois ou \acourt{V}{10, exercice 20}.) \end{proof} Comme
$f$ est quasi-compact puisqu'affine, le lemme est prouvé.  \end{proof}

Dans un second temps, rappelons la construction de la topologie constructible sur $X$ (cf.
\egalong{IV}{3}{}).

\begin{remarque2} Pour notre propos, on ne l'utilisera en fait que pour le schéma affine
$X'$.\end{remarque2}

On construit l'espace topologique $X^{\co}$ dont l'ensemble sous-jacent coïncide avec l'ensemble
sous-jacent $|X|$  de $X$ mais dont les ouverts (resp. fermés) sont les parties ind (resp.
pro)-constructibles, à savoir les réunions (resp. intersections) de parties constructibles.  Comme
$X$ est cohérent, $X$ est un espace topologique compact, totalement discontinu
(\egalong{IV}{1}{1.9.15}). De plus, la cohérence de $X$ entraîne que les parties constructibles sont
alors les réunions finies d'intersection $U\cap(X-V)$ avec $U,V$ ouverts quasi-compacts
(\egalong{III}{1}{9.1.3} et \egalong{IV}{1}{1.2.7}). Le complémentaire de $U\cap(X-V)$ étant
$(X-U)\cup V$, il est donc également ouvert dans $X^{\co}$ de sorte que $X$ admet une base d'ouverts
compacts.

L'identité de $|X|$ induit une application continue $X^{\co}\ra X$ puisqu'un fermé est
pro-constructible. Si $X=\Spec(A)$ est affine, $X^{\co}$ est naturellement homéomorphe au spectre
d'une certaine $A$-algèbre $T(A)$ pour un certain endo-foncteur $T$ de la catégorie des anneaux
(\cite{Oli67}, proposition 5). Cet homéomorphisme est compatible à la localisation de sorte que ces
structures schématiques se recollent munissant $X^{\co}$ d'une structure naturelle de $X$-schéma
relativement affine compatible avec l'application continue (identique !) $X^{\co}\ra X$. Si $x\in
|X|$, on a $\OO_{x}=k(x)$.


On définit alors $$f:X^c\ra X$$ comme le composé $$f:X^c=(X')^{\co}\ra X'\ra X.$$

Par construction, $X^c$ est compact -en particulier cohérent- (réduit), totalement discontinu et
admet une base de voisinages ouverts-compacts (qui sont donc ouverts-fermés puisque $X^c$ est
compact donc topologiquement séparé).  Ses corps résiduels sont algébriquement clos et $f$ est
quasi-compact\footnote{On appliquera ici cette construction à un ouvert (quasi-compact) d'un schéma
affine, donc à un schéma séparé de sorte que $f$ sera même affine dans ce cas.} et surjective (comme
composé de morphismes quasi-compacts) .

Rappelons que le faisceau vide sur un espace topologique est le faisceau associé au préfaisceau de
valeur constante $\varnothing$.  L'ensemble de ses sections sur tout ouvert non vide est
$\varnothing$ et est réduit à un point sur l'ouvert vide.

\begin{lemme2}\label{20-co-sec} Tout morphisme étale $f:Y\ra X^c$ est Zariski localement trivial.
En particulier, le morphisme canonique de topos $(\epsilon^{-1},\epsilon_*): X^c_\et\ra
X^c_{\textup{Zar}}$ est une équivalence d'inverse $(\epsilon^{-1},\epsilon_*)$. Tout faisceau ayant
des sections localement a des sections globales. De plus, tout torseur sur $X^c$ est trivial et
toute gerbe est neutre.\end{lemme2}

\begin{proof} Soit $y\in Y$ d'image $x\in X^c$. Comme $f$ est quasi-fini et $k(x)$ algébriquement
clos, l'inclusion $k(x)\hookrightarrow k(y)$ est une égalité. Le morphisme composé
$$\Spec(\OO_x)=\Spec(k(x))=\Spec(k(y))\ra Y$$ se prolonge au voisinage de $f$ en une section locale
de $f$ ce qui prouve le premier point.

Soit $\F$ un faisceau (Zariski ou étale, c'est la même chose) sur $X^c$. Il est donc Zariski
localement trivial. Comme $X^c$ est compact, on peut trouver un recouvrement fini par des ouverts
compacts $U_i$ sur lesquels $\F$ a une section. On montre par récurrence sur le nombre d'ouverts
qu'on peut raffiner ce recouvrement en un recouvrement fini $V_j$ par des ouverts compacts
disjoints. Comme pour tout $j$ il existe $i$ tel que $V_j\subset U_i$, le faisceau $\F$ a des
sections locales sur chaque $V_j$. Ces ouverts étant disjoints, ces sections se recollent en une
section globale. Le reste suit car tout torseur (resp. toute gerbe) sur $X^c$ a des sections
localement.\end{proof}

La preuve du lemme d'effacement est alors facile. Soit $\gamma$ une gerbe ind-finie sur $X$. Le
foncteur d'adjonction $$\gamma\ra f_*f^* \gamma$$ est fidèle car $f$ est surjectif. La gerbe
$f^*\gamma$ (\cite{CNA@Giraud}, III.2.1.5.6) est neutre et ind-finie (\ref{20-co-sec}) de sorte
qu'elle est équivalente à $\underline{\Tors}(X^c,G^c)$ pour un ind-groupe $G^c$ convenable. Par
ailleurs, comme on a déjà vu (\cite{CNA@Giraud}, V.3.1.9.1), le faisceau des sous-gerbes maximales
$\pi_0(f_*\underline{\Tors}(X^c,G^c))$ s'identifie à $R^1f_*G^c$.

\begin{lemme2} Le faisceau $R^1f_*G^c$ est trivial.\end{lemme2} \begin{proof} Soit $U\ra X$ un
morphisme étale. Comme dans la preuve de \ref{20-co-sec}, $X^c_U\ra X^c$ est étale et donc un
isomorphisme local pour la topologie de Zariski car les corps résiduels de $X^c$ sont algébriquement
clos. On en déduit que les topos étale et Zariski de $X^c_U$ sont équivalents. Tout $G^c$-torseur
étale sur $X_U^c$ provient donc d'un torseur Zariski.  Comme $X^c_U\ra X^c$ est un isomorphisme
local, $X_U^c$ a une base d'ouverts-fermés et donc est séparé. Comme $X^c_U$ est quasi-compact
(puisque $f$ est quasi-compact et $U$ quasi-compact comme $X$), l'espace topologique sous-jacent de
$X^c_U$ est de plus compact. On déduit comme dans \ref{20-co-sec} que tout torseur sur $X_U^c$ est
trivial $H^1(X_U^c,G^c)=\{*\}$. En passant à la limite (on n'utilise pas ici la cohérence de $f$),
on trouve que les fibres de $R^1f_*G_c$ sont triviales.\end{proof}


D'après le lemme, $\pi_0(f_*\underline{\Tors}(X^c,G^c))$ est le faisceau ponctuel ce qui assure que
$f_*\underline{\Tors}(X^c,G^c)$ est une gerbe. Comme elle a une section, elle est neutre de groupe
$G=f_*G^c$ et s'identifie donc à $\underline{\Tors}(X,G)$. Mais $f$ est quasi-compact de sorte que
$G$ est ind-fini comme $G^c$ \sga{4}{IX}{1.6}. La preuve du lemme d'effacement est complète.


\section{Appendice 2: théorème de changement de base propre  d'Artin-Grothendieck pour les champs
ind-finis sur des schémas non noethériens}

On va prouver l'énoncé suivant

\begin{theoreme}\label{20-th-cgt-base-propre}Considérons un
diagramme cartésien $$\xymatrix{X'\ar[r]^{g'}\ar[d]_{f'}\ar@{}[rd]|\square&X\ar[d]^f\\
Y'\ar[r]^g&Y}$$ avec $f$ propre. Alors, pour tout champ ind-fini $\C$, la flèche de changement de
base $g^*f_*\C\ra f'_*g'^*\C$ est une équivalence.
\end{theoreme}

\begin{proof} Notons que ce théorème est connu dans le cas discret \sga{4}{XII}{5.1} et si $Y$ est
localement noethérien (\cite{CNA@Giraud}, VII.2.2.2). La preuve qui suit est une adaptation de la
preuve de ce dernier énoncé.

\textit{Preuve du théorème : } D'après \cite{CNA@Giraud}, VII.2.2.5, il suffit de prouver l'énoncé
suivant : soit $X$ propre sur $S$ local hensélien et $i:X_0\hookrightarrow X$ l'immersion de la
fibre fermée. Alors, pour tout champ ind-fini $\C$ sur $X$, la flèche \begin{equation*} \gamma:
H^0(X,\C)\ra H^0(X_0,i^*\C) \end{equation*} est une équivalence. Notons que $X/S$ étant propre, il
est cohérent donc $X$ est cohérent comme $S$ \sga{4}{VI}{2.5}.  Ainsi, $i$ est un morphisme cohérent
de schémas cohérents. Si $\C$ est discret, le théorème est une conséquence immédiate du théorème de
changement de base propre pour les faisceaux d'ensembles \sga{4}{XII}{5.1 (i)}.  On en déduit que
$\gamma$ est pleinement fidèle. D'après \ref{20-prop-red-champ}, il suffit de montrer que pour tout
morphisme fini $X'\ra X$ (induisant une immersion fermée $X'_0\hookrightarrow X'$) et tout groupe
fini $G$, la flèche \begin{equation}\label{20-rest-ferm} \gamma: \Tors(X',G)\ra \Tors(X'_0,G)
\end{equation} est une équivalence (et en fait est essentiellement surjective puisqu'on sait déjà
qu'elle est pleinement fidèle). On applique alors \sga{4}{XII}{5.5 (ii)} au morphisme propre $X'\ra
S$ pour conclure.\end{proof}


\section{Appendice 3: sorites sur les gerbes}

On montre que toute gerbe ind-finie sur $X$ noethérien est limite inductive filtrante de ses
sous-gerbes constructibles (comparer avec \egalong{IV}{3}{IX.2.9} et \egalong{IV}{3}{IX.2.2}).

\subsection{Image d'un morphisme de champs} Soit $\varphi:\C\ra\C'$ un morphisme (cartésien) de
champs sur $X_\et$.  On définit l'image $\Im(\varphi)=\varphi(\C)$ comme la catégorie ayant pour
objets ceux de $\C$ et telles que
$\Hom_{\varphi(\C)}(g_1,g_2)=\Hom_{\C'}(\varphi(g_1),\varphi(g_2))$, la structure de catégorie
fibrée étant déduites des structures (compatibles) de $\C$ et $\C'$. Notons que $\Im(\varphi)$ est
naturellement équivalente à la sous-catégorie pleine  $\C'_\varphi$ de $\C'$ dont les objets sont
les images des objets de $\C$. On identifiera $\varphi(\C)$ et $\C'_\varphi$.

\begin{lemme2}\label{20-lemme-im} Soit $\bar x$ un point géométrique de $X$ et $f:W\ra X$ un
morphisme de schémas. Alors, \begin{itemize} \item la flèche naturelle $\Im(\varphi)_{\bar
x}\ra\Im(\varphi_{\bar x})$ est une équivalence ; \item on a une équivalence de champs canoniques
$f^*\Im(\varphi)\iso\Im(f^*\varphi)$.  \end{itemize} \end{lemme2}

\begin{proof} Les objets de $\Im(\varphi)_{\bar x}\ra\Im(\varphi_{\bar x})$ coïncident avec ceux de
$\C_{\bar x}$ et la flèche naturelle est simplement l'identité. Construisons l'inverse de la flèche.
Soit alors $a,b\in\bar C_{\bar x}$ et $\psi\in\Hom(\phi(a),\phi(b))$ qui provient de
$\Psi_S\in\Hom_S(\phi(\alpha),\phi(\beta))$. Mais $\Psi_S$ peut-être vu comme une flèche de
$\Im(\varphi)(S)$ : on prend son germe en $\bar x$ pour définir l'inverse (qui ne dépend pas des
choix). On vérifie que ceci définit l'inverse cherché.

Passons au second point et définissons la flèche. Par adjonction, on doit définir une flèche
(cartésienne) $$\Im(\varphi)\ra f_*\Im(f^*\phi).$$


Soit $S\ra X$ étale. Les objets du membre de gauche sont les objets de $\C(S)$ tandis que ceux de
droites sont ceux de $f^*\C(f^{-1}(S))=f_*f^*\C(S)$. La flèche d'adjonction $\C\ra f_*f^*\C$ permet
alors de définir  la flèche $x\mapsto f^*(x)$ cherchée au niveau des objets. Soient alors $x,y$ des
objets de $\C(S)$ et $$g\in\Hom_S(\varphi(x),\varphi(y))=\Hom_{\Im(\varphi)(S)}(x,y)$$ C'est donc
une section sur $S$ de $\underline{\Hom}(\varphi(x),\varphi(y))$ qui fournit (par image
inverse) une section sur $f^{-1}(S)$ de $\underline{\Hom}(f^*\varphi(x),f^*\varphi(y))$
(\cite[II.3.2.8.1 (4)]{CNA@Giraud}, donc une flèche de
$$\Hom_{f^{-1}(S)}(f^*\varphi(x),f^*\varphi(y))=\Hom_{f^{-1}(S)}((\varphi(f^*x),\varphi(f^*y))=
\Hom_{f_*\Im(f^*\varphi)(S)}(x,y).$$ Le foncteur ainsi défini est visiblement cartésien (comme
$\varphi$).  Le premier point assure que les fibres de ce foncteur sont des équivalences, ce qui
achève de prouver le lemme.\end{proof}

\subsection{Groupoïdes libres}Soit $\Gamma=\xymatrix{E\ar@<1ex>[r]^s\ar@<-1ex>[r]_b&V}$ un graphe
orienté ($E$ est l'ensemble des arêtes, $V$ l'ensemble des sommets, $b,s$ les applications « but,
source »). On associe (voir \cite{Ber95}) le groupoïde libre $L(\Gamma)$ qu'on peut décrire comme
suit. Soit $E^{\pm}$ l'ensemble $E^\pm=\{e^{\pm},e\in E\}$ union disjointe de deux copies de $E$ :
ses objets sont les sommets et les morphismes entre $v,v'\in V$ sont les mots (réduits)
$$e_1^{\pm}\cdots e_n^{\pm} \text{ avec } b(e_i)=s(e_{i+1}) (i=1,\cdots, n-1),  s(e_1)=v,
b(e_n)=v'.$$

\begin{remarque2} Il est bien connu que  $L(\Gamma)$ est le groupoïde fondamental
$\Pi_1(\Gamma_\RR)$ de la réalisation géométrique $\Gamma_\RR$ de $\Gamma$.\end{remarque2}

Par construction, les foncteurs de $L(\Gamma)$ dans un groupoïde $G$ s'identifient naturellement aux
familles $$(g_v)\in \Ob(G)^V, (\gamma_e)\in\mathrm{Fl}(G)^E\text{ telles que }\gamma_e\in
\Hom_G(g_{s(e)},g_{b(e)}).$$ Si on préfère, $L$ est l'adjoint à gauche du foncteur d'oubli
$\text{Groupoïdes}\ra\text{Graphes}.$

La construction se globalise de la manière suivante. Considérons un diagramme de $X$-schémas étales

$$\Gamma_X:\xymatrix{E\ar@<1ex>[rr]^s\ar@<-1ex>[rr]_b\ar[rd]&&V\ar[ld]\\&X}$$

Par fonctorialité de la construction $L$, on définit un préchamp sur $X_\et$ par la formule
$S\mapsto L(\Gamma_X(S))$ dont le champ associé est noté $L(\Gamma_X)$ : c'est le groupoïde libre
engendré par le graphe orienté. Par construction $L(\Gamma_X)$ a des sections locales si et
seulement si $V\ra X$ est surjectif. Les sections locales sont localement isomorphes si et seulement
si pour tout point géométrique $\bar x\ra X$, le graphe  $\xymatrix{E_{\bar
x}\ar@<1ex>[r]\ar@<-1ex>[r]&V_{\bar x}}$ est connexe. Par construction, on dispose de deux sections
tautologiques $$g\in \Ob(\Gamma_X)(V),\gamma\in\Hom_{(\Gamma_X)(E)}(s^*g,b^*g)$$  définies par
l'identité de $V$ et de $E$ respectivement. Soit $\G$ un groupoïde sur $X_\et$. On a alors la
propriété d'adjonction suivante : la flèche qui à un foncteur \emph{cartésien}
$$\varphi:L(\Gamma_X)\ra\G$$ associe
$$\varphi(g)\in\G(V),\varphi(\gamma)\in\Hom_{\G(E)}(s^*\varphi(g),b^*\varphi(g))$$ est bijective.



\subsection{Constructibilité de sous-gerbes} Considérons un hyper-recouvrement\footnote{La
terminologie est abusive : manque la section diagonale $V\ra E$ pour avoir un hyper-recouvrement
(tronqué).} de $X$, à savoir un diagramme de $X$-schémas étales (de type fini)

$$H_X:\xymatrix{E\ar@<1ex>[rr]^s\ar@<-1ex>[rr]_b\ar[rd]&&V\ar[ld]\\&X}$$ où les flèches $V\ra X$ et
$(s,b):E\ra V\times_XV$ sont surjectives. Pour tout point géométrique $\bar x\ra X$, le graphe
$H_{\bar x}$ est connexe de sorte que $L(H_X)$ est une gerbe. Soit $$g\in
\G(V),\gamma\in\Hom_{\G(E)}(s^*\varphi(g),b^*\varphi(g))$$ définissant un morphisme cartésien
$\varphi :L(H_X)\ra\G$.



\begin{definition2} Une gerbe $\G$ sur $X_\et$ est dite constructible si pour toute section locale
$\sigma\in\G(S), S\ra\X$ ouvert étale, le faisceau en groupes $\underline{\Aut}(\sigma)$ sur $S_\et$
est constructible.\end{definition2}


\begin{lemme2} Avec les notations précédentes, supposons $\G$ ind-finie. Alors, l'image
$I=\Im(\varphi)$ est constructible.\end{lemme2}

\begin{proof} Comme la formation de l'image et de $L$ commutent à l'image inverse, on peut procéder
par récurrence noethérienne. Il suffit donc de prouver que $I$ est constructible sur un ouvert non
vide de $X$. La constructibilité se testant après n'importe quel changement de base surjectif
localement de présentation finie,\sga{4}{IX}{2.8}, on peut supposer que $V,E$ sont des
\emph{revêtements} étales de $X$ complètement décomposés, autrement dit que $H_X$ est un graphe
constant $\Gamma$ de sorte que $\sigma$ s'identifie à $n=\card(V)$ sections $\sigma_i\in\G(X)$ deux
à deux isomorphes. Ainsi, $L(H_X)$ est le champ en groupoïdes constant $L(\Gamma)_X$. Soit $N$ le
noyau du morphisme de groupoïdes (abstraits) $$N=\Ker(L(\Gamma)\ra \G(X)).$$ Par adjonction,
$L(H_X)=L(\Gamma)_X\ra \G$ se factorise à travers la projection $L(\Gamma)_X\ra (L(\Gamma)/N)_X$.
Mais, comme $L(\Gamma)/N$ est un groupoïde \emph{fini}, il est constructible de sorte que $I$ l'est
aussi, comme quotient de $(L(\Gamma)/N)_X$.\end{proof}

\begin{proposition2}\label{20prop-constr} Soit $\pi:X\ra Y$ un morphisme de schémas noethériens,
$\G$ une gerbe ind-finie sur $Y_\et$ et $\sigma\in H^0(X,\pi^*\G)$. Il existe  une sous-gerbe
constructible $\G_1$ de $\G$ telle que $\sigma\in H^0(X,\pi^*\G_1)\subset H^0(X,\pi^*\G)$.
\end{proposition2}

\begin{proof} La formule $\pi^*\G_{x}=\G_{\pi(x)}$ assure que localement $\sigma$ provient par image
inverse d'une section locale de $\G$. Comme $Y$ est quasi-compact, on peut trouver $V\ra Y$ étale
(surjectif de type fini) et $\tau\in \G(V)$ telles que $\pi^*\tau$ et $s$ coïncident localement sur
$X$. En considérant $E\ra V\times_YV$ étale convenable définissant un isomorphisme $\gamma:
s^*\tau\ra b^*\tau$, il suffit d'après le lemme précédent de poser
$\G_1=\Im(L(V,E)\xrightarrow{(\tau,\gamma)}\G)$.\end{proof}

\begin{remarque2} On contourne ici l'absence de sorites sur les limites inductives. L'énoncé devrait
être en deux parties : d'abord qu'une gerbe ind-finie sur un schéma noethérien est limite inductive
filtrante de ses sous gerbes constructibles, ce qui est pour l'essentiel le contenu du lemme
précédent, ensuite que sur un schéma cohérent, le foncteur sections globales commute aux limites
inductives filtrantes.\end{remarque2}



\ifx\danslelivre\undefined \bibliography{../fichiers-configuration-latex/bibliographie}
\bibliographystyle{../fichiers-configuration-latex/style-GTG} \end{document} \fi

